\chapter{Elementos básicos de cálculo multivariado}

Presentamos en este capítulo algunos de los elementos teóricos básicos que se requieren para abordar 
los problemas planteados en los capítulos posteriores.
\vskip 0.1cm
Usaremos el término \emph{conjunto} para referirnos a una colección bien definida de objetos. 
Una \emph{función} $f$ de un conjunto $A$ 
en un conjunto $B$, es una regla que asigna a cada elemento de $A$ un único elemento de $B.$
Seguramente el lector está familiarizado con el conjunto de números reales $\mathbb{R}.$
En estos apuntes consideraremos funciones cuyo dominio es el espacio el espacio $n-$dimensional $\mathbb{R}^n,$
o un subconjunto de este.

\section{Vectores}
En \cite{marsden_calculo_2013} se define "vector como un segmento de recta dirigido que comienza en el origen,
esto es, un segmento de recta con magnitud y dirección específicos con punto inicial en el origen".\\
Los puntos en $\mathbb{R}^n$ se presentan en la forma $P=(a_1,a_2,\cdots, a_n)$ 
y el vector que va desde el origen hasta el punto $P$ se denotará por $\langle a_1,a_2,\cdots, a_n\rangle.$ 
Particularmente a los vectores $\vi=\langle 1,0,0 \rangle,$ $\vj=\langle 0,1,0 \rangle,$ $\vk=\langle 0,0,1 \rangle$ los denominaremos \emph{vectores coordenados} en $\mathbb{R}^3.$ 
\vskip0.2cm
Para dos vectores $\ve{u}=\langle a_1,a_2,\cdots, a_n\ra$ y $\ve{v}=\langle b_1,b_2,\cdots, b_n\ra$
el \emph{producto punto} se denota por $\ve{u}\cdot \ve{v} = a_1b_1+a_2b_2+\cdots+a_nb_n$; 
la \emph{norma} $\ve{u}$ se denota por $\Vert \ve{u} \Vert$ y se define
por  $\Vert \ve{u}\Vert=\sqrt{a_1^2+a_2^2+\cdots+a_n^2};$ nótese que
$\ve{u}\cdot \ve{u}=\Vert \ve{u}\Vert^2.$
El \'{a}ngulo $\theta $ entre los vectores $\ve{u}$ y $\ve{v}$ satisface la relación
$\ve{u} \cdot \ve{v}= \left\|\ve{u}\right\| \left\| \ve{v}\right\| \cos \theta .$
\vskip0.2cm
El \emph{producto cruz} entre dos vectores  $\ve{u}=\langle a_{1},a_{2},a_{3}\ra$ y 
$\ve{v}=\langle b_{1},b_{2},b_{3}\ra$ en $\mathbb{R}^3$ es el vector 
$$\ve{u}\times \ve{v}=\langle a_{2}b_{3}-a_{3}b_{2},a_{3}b_{1}-a_{1}b_{3},a_{1}b_{2}-a_{2}b_{1}\ra,$$
que es ortogonal tanto a $\ve{u}$ como a $\ve{v}.$ 
Adem\'{a}s $\left\Vert \ve{u}\times \ve{v}\right\Vert =\left\Vert \ve{u}\right\Vert\left\Vert \ve{v}\right\Vert\sin\theta;$ estos resultados se pueden consultar en detalle en \textcite{apostol_calculus_1999}.

La siguiente es una conocida forma de evaluar este producto,
$$\ve{u}\times \ve{v}=\left| \begin{array}{ccc} \vi & \vj & \vk \\ a_1 &a_2&a_3 \\b_1 &b_2&b_3 \\ \end{array} \right|  $$
El producto cruz es útil en el desarrollo de algunos de los siguientes resultados.
\subsection*{Proyección ortogonal}
%
La proyecci\'{o}n ortogonal del vector $\ve{v}$ sobre el vector 
$\ve{u},$ se denota por ${\rm Proy}_{\ve{u}}\ve{v}$
y es un múltiplo escalar del vector $\ve{u}.$ Es decir 
${\rm Proy}_{\ve{u}}\ve{v}=k\ve{u}.$
De la figura se observa que 
$$\cos \theta =\frac{\left\| k\ve{u}\right\|}{\left\|
\ve{v}\right\|}=\frac{\left| k\right| \left\| \ve{u} \right\|}{\left\|\ve{v}\right\|}.$$


El valor $\left| k\right| $ ser\'{a} positivo si 
${\rm Proy}_{\ve{u}}\ve{v}$ posee el mismo sentido de $\ve{u},$ de lo
contrario, ser\'{a} negativo. Al igualar las expresiones para $\cos \theta $ y considerando el caso $k>0,$
se concluye que $k=\frac{\ve{u}\cdot \ve{v}}{\left\| \ve{u}\right\| ^{2}},$ entonces 
la proyección del vector $\ve{v}$ sobre $\ve{u}$ está dada por
$${\rm Proy}_{\ve{u}}\ve{v}=\frac{\ve{u}\cdot\ve{v}}{\left\| \ve{u}\right\|^{2}}\ve{u}$$
En el gr\'{a}fico la línea discontinua corresponde al vector 
$\ve{v}-{\rm Proy}_{\ve{u}}\ve{v}$. Aunque nos apoyamos en un gráfico bidimensional, este resultado se puede presentar en espacios 
de dimensión superior.

\vskip0.2cm
\noindent
\begin{minipage}{\textwidth}
	\centering
	\captionof{figure}{Proyección ortogonal}
	\label{fig:proy-ortogonal}
	\begin{minipage}{0.30\linewidth}
		\centering
		\includegraphics[width=\linewidth]{Fig/1.pdf}
	\end{minipage}
	
	\vspace{0.3em}
	\parbox{0.95\linewidth}{\centering\fontsize{8pt}{10pt}\selectfont\textbf{Fuente:} elaboración propia.}
\end{minipage}
\vskip0.2cm


\begin{prop}
Para todo $\ve{v}\in \mathbb{R}^n$, el vector $\ve{v}-{\rm Proy}_{\ve{u}}\ve{v}$ es ortogonal
tanto a $\ve{u}$ como a ${\rm Proy}_{\ve{u}}\ve{v}.$
\end{prop}
\begin{proof}
Veamos que $\ve{v}-{\rm Proy}_{\ve{u}}\ve{v}$ es ortogonal a $\ve{u}$
\begin{eqnarray*}
\ve{u} \cdot \left(\ve{v}-{\rm Proy}_{\ve{u}}\ve{v}\right) &=& \ve{u} \cdot \left(\ve{v} - \frac{\ve{u}\cdot \ve{v}}{\Vert\ve{u}\Vert^2}\ve{u}\right) \\ 
&=& \ve{u}\cdot\ve{v}-\left (\frac{\ve{u}\cdot\ve{v}}{\Vert\ve{u}\Vert^2}\right)\ve{u}\cdot\ve{u} \\
&=& \ve{u}\cdot\ve{v}-\left (\frac{\ve{u}\cdot\ve{v}}{\Vert\ve{u}\Vert^2}\right) \Vert\ve{u}\Vert^2 \\
&=& \ve{u}\cdot\ve{v}-\ve{u}\cdot\ve{v}=0.
\end{eqnarray*}

Como ${\rm Proy}_{\ve{u}}\ve{v}$ es múltiplo escalar de $\ve{u},$ con el resultado anterior se sigue que
$\ve{v}-{\rm Proy}_{\ve{u}}\ve{v}$ es ortogonal a ${\rm Proy}_{\ve{u}}\ve{v}.$
\end{proof}
\subsection*{Área de paralelogramo}
El área $A$, del paralelogramo cuyas aristas son los vectores $\ve{u}$ y $\ve{v}$
está dada por $\Vert \ve{u} \Vert h = \Vert \ve{u} \Vert \Vert \ve{v} \Vert \sin \theta . $
De la relación $\Vert \ve{u} \times \ve{v} \Vert =\Vert \ve{u} \Vert \Vert \ve{v} \Vert \sin \theta$
se sigue que el área del paralelogramo es
$A=\Vert \ve{u} \times \ve{v} \Vert.$


\begin{figure}
    \centering
    \caption{Área de un paralelogramo y volumen de un paralelepípedo}
    \label{paralelogramo_paralelipedo}
    \includegraphics[width=0.38\linewidth]{Fig/2.pdf}
    \hspace{1em}
    \includegraphics[width=0.38\linewidth]{Fig/3.pdf}  
    \parbox{0.95\linewidth}{\centering\fontsize{9pt}{10pt}\selectfont\textbf{Fuente:} elaboración propia.}
\end{figure}

\subsection*{Volumen de un paralelepípedo}
El volumen del paralelepípedo cuyas aristas son los vectores $\ve{u},$ $\ve{v}$ y $\ve{w},$
se determina con el siguiente procedimiento: 
el área de la base del paralelepípedo está dada por $\Vert \ve{u} \times \ve{v}  \Vert$,
la altura $h$ del paralelepípedo corresponde a la norma de la proyección del vector $\ve{w}$
sobre el vector $\ve{u} \times \ve{v}.$
Es decir,
$ h  = \Vert {\rm Proy}_{\ve{u} \times \ve{v}} \ve{w} \Vert =\left\Vert \frac{\ve{w} \cdot (\ve{u} \times \ve{v})}{\Vert\ve{u} \times \ve{v} \Vert^2 }\,(\ve{u} \times \ve{v} )\right \Vert = \frac{|\ve{w} \cdot (\ve{u} \times \ve{v})|}{\Vert\ve{u} \times \ve{v} \Vert }$. 
Por lo anterior, el volumen $V$ del paralelepípedo es $h \Vert\ve{u} \times \ve{v} \Vert$
entonces, 
\begin{equation} \label{preliminares4}
V= |\ve{w} \cdot (\ve{u} \times \ve{v}) |
\end{equation}

\section{Funciones en el espacio} \label{funcionesenelespacio}
Para $a\in \mathbb{R}^{n},$ al conjunto $B_{r}\left(a\right) =\{x\in \mathbb{R}^{n}|\left\Vert x-a\right\Vert <r\}$ 
lo denominaremos la bola con centro en $a$ y radio $r>0.$ \index{Bola} 
Diremos que un conjunto $U\subseteq \mathbb{R}^{n}$ es \emph{abierto} si para todo $u\in U$ 
existe $r>0,$ tal que, $B_{r}(u) \subset U.$  \index{Conjunto!abierto}
Un sinónimo de conjunto abierto es región abierta. \index{Región abierta} 
Un conjunto $B$ se dice \emph{cerrado} si su complemento es un conjunto abierto. \index{Conjunto!cerrado}
\begin{teorema}
$B_{r}(a)$ es un conjunto abierto.
\end{teorema}
La demostración de este resultado se encuentra en \textcite{marsden_calculo_2013}.
\begin{ejemplo}
El conjunto $A=\{(x,y) \in \mathbb{R}^{2}|x>0\}$ es abierto pues
para todo $a \in A$ se satisface la definición dada.
\vskip0.2cm
El conjunto $B=\{(x,y) \in \mathbb{R}^{2}|xy\geq 0\}$ no es abierto pues
para $\left( 0,0\right) \in B$ no es posible determinar
$r>0$ tal que $B_{r}\left( \left( 0,0\right) \right) \subset B.$
\end{ejemplo}
De manera general abordaremos funciones de la forma
\begin{equation}
f: A\subseteq \mathbb{R}^n \to B\subseteq \mathbb{R}^m
\end{equation}
El conjunto de puntos para los cuales la función $f$ está definida se denomina \emph{dominio} de $f$ 
y el conjunto de puntos de $\mathbb{R}^m$ que son imagen de algún punto del dominio se denomina \emph{rango}
de $f$. De acuerdo con \textcite{apostol_calculus_1999}, a las funciones definidas cuando $n=m=1,$ se denominan \emph{funciones reales,} las cuales se estudian en los cursos de \emph{cálculo diferencial} y \emph{cálculo integral.}
Si $n>1$ y $m=1$ obtendremos \emph{campos escalares};  \index{Campos escalares} \index{Campos vectoriales}
llamaremos \emph{funciones vectoriales} a las obtenidas cuando $n=1$ y $m>1$; mientras que los \emph{campos vectoriales} son funciones obtenidas cuando $n>1$ y $m>1.$ \index{Funciones vectoriales}
\vskip0.2cm 
Para presentar funciones vectoriales, campos vectoriales y en general para vectores, usaremos la notación $\ve{r}(t),$ ${\bf r}(t)$ o ${\bf F}(x,y,z)$ indistintamente y según el caso. 
Por ejemplo:

\begin{itemize}
\item $\overrightarrow{r}(t) =\left\langle \ln t,\frac{1}{t},t^{2}\right\rangle ,$ es 
una funci\'{o}n vectorial en $\mathbb{R}^{3}$ con dominio $(0,\infty ).$ 
% El rango es el conjunto $B={\rm Rango}f=\{(x,y,z)\in\mathbb{R}^{3}/ y\neq 0, z \geq 0\}.$
\item El campo escalar definido en $\mathbb{R}^{2}$ por $f(x,y) =\sqrt{9-x^{2}}+\sqrt{y^{2}-4}$ posee
dominio, ${\rm dom}f=\{(x,y) \in \mathbb{R}^{2}/\left\vert x\right\vert \leq 3,\left\vert y\right\vert \geq 2\}$ y el rango de $f$ es el conjunto $[0,\infty).$
\item $\overrightarrow{F}(x,y,z) =\left\langle y,\frac{x}{x-y},e^{-xz}\right\rangle$ 
es un campo vectorial definido en $\mathbb{R}^{3},$ su dominio es el conjunto 
${\rm dom}f=\{(x,y,z) \in \mathbb{R}^{3}/x\neq y\}.$
\end{itemize}
Supondremos que el lector ha estudiado los conceptos de límite, continuidad, diferenciabilidad e integrabilidad
en funciones reles, de tal manera que extenderemos estos elementos a los otros tipos de funciones. 
\vskip0.2cm 
En cuanto al dominio de la funci\'{o}n $f(x,y) =\frac{x^{4}-y^{4}}{x^{2}+y^{2}},$
${\rm Dom}f=\{(x,y) \in \mathbb{R}^{2}/(x,y) \neq \left( 0,0\right)\}$
a pesar que $\left( 0,0\right) $ no es un punto del dominio de $f,$
es natural preguntarnos: ¿qué sucede si nos acercamos al origen ? En la siguiente sección podremos 
responder a este interrogante.

\subsection*{Límites}

\begin{defi} \label{preliminares5}
Si {\small $\ve{F}: A \subseteq\mathbb{R}^n \to  \mathbb{R}^m$} es una función con dominio $A,$
se dice que $\ve{L}$ es el límite de $\ve{F}(x_{1},x_{2},\cdots ,x_{n})$ cuando 
${\bf x}=\left(x_{1},x_{2},\cdots ,x_{n}\right)$ tiende a ${\bf a}= \left( a_{1},a_{2},\cdots,a_{n}\right) ,$ 
lo cual se denota como {\small $ \lim_{{\bf x}\to {\bf a}}\ve{F}({\bf x})=\ve{L}=\langle l_{1},l_{2},\cdots ,l_{m}\ra,$}
si para todo {\small $\varepsilon >0$} existe $\delta >0$ tal que {\small $\left\Vert \ve{F}\left( x_{1},x_{2},\cdots ,x_{n}\right) -\ve{L}\right\Vert <\varepsilon,$}
siempre que $\left\Vert {\bf x}- {\bf a} \right\Vert <\delta .$
\end{defi} 

\begin{itemize}[leftmargin=0cm,itemindent=0.4cm,labelwidth=\itemindent,labelsep=0cm,align=left]
\item Usando la definición, demostrar que $\lim\limits_{(x,y) \rightarrow \left( 1,1\right) }\left\langle
x^{2},x+y\right\rangle =\left\langle 1,2\right\rangle .$
\vskip0.1cm

Observemos que  para todo $\varepsilon >0$ existe $\delta >0$ tal que si 
$\left\Vert(x,y) -\left( 1,1\right) \right\Vert <\delta $
entonces $\left\Vert \left( x^{2},x+y\right) -\left( 1,2\right) \right\Vert<\varepsilon ;$ esto se evidencia verificando las siguientes desigualdades:
\begin{eqnarray}\label{preliminares}
\left\vert x-1\right\vert &=&\sqrt{(x-1) ^{2}}\leq \sqrt{\left(
x-1\right) ^{2}+(y-1) ^{2}} \nonumber \\ &=&\left\Vert \left(x,y\right) -\left( 1,1\right) \right\Vert <\delta
\end{eqnarray}
De esta relaci\'{o}n se obtiene
\begin{multline}\label{preliminares1}
|x-1|<\delta \rightarrow -\delta <x-1<\delta
\rightarrow 2-\delta <x+1<\delta  \\
\Rightarrow \left\vert x+1\right\vert<2-\delta 
\end{multline}
De manera an\'{a}loga,
\begin{eqnarray}\label{preliminares2}
\left\vert y-1\right\vert &=& \sqrt{(y-1) ^{2}}\leq \sqrt{\left(x-1\right) ^{2}+(y-1)  ^{2}} \nonumber \\ &=& \left\Vert \left(x,y\right) -\left( 1,1\right) \right\Vert <\delta 
\end{eqnarray}
Utilizando los resultados (\ref{preliminares}), (\ref{preliminares1}) y (\ref{preliminares2}), se sigue que
\begin{eqnarray*}
\left\Vert \langle x^{2},x+y\rangle -\langle 1,2\rangle \right\Vert &=& \sqrt{\left( x^{2}-1\right) ^{2}+\left( x+y-2\right) ^{2}} \\
&\leq & \sqrt{\left( x^{2}-1\right) ^{2}} +\sqrt{\left(x+y-2\right) ^{2}} \\
&=& \left\vert x^{2}-1\right\vert  +\left\vert x+y-2\right\vert \\
&=& \left\vert x-1\right\vert  \left\vert x+1\right\vert +\left\vert (x-1) +(y-1) \right\vert \\
&<& \delta (2-\delta) +2\delta = \delta (4-\delta). 
\end{eqnarray*}
Tomando $\varepsilon =\delta \left( 4-\delta \right) $, se obtiene que 
$\delta =2-\sqrt{4-\varepsilon };$ con $0<\varepsilon < 4$ la definici\'{o}n se satisface.

\vskip0.3cm

\item Mostrar que $\lim\limits_{t\rightarrow 1}\left\langle \frac{t^{2}-1}{t-1},\frac{t+1}{t},t^{2}\right\rangle =\left\langle 2,2,1\right\rangle .$ Veamos que para todo $\varepsilon >0$ existe 
$\delta >0$ tal que si $\left\vert t-1\right\vert <\delta, $ 
entonces $\left\Vert \left\langle \frac{t^{2}-1}{t-1},\frac{t+1}{t},t^{2}\right\rangle -\left\langle 2,2,1\right\rangle \right\Vert<\varepsilon .$
De la relación
\begin{eqnarray*}
\left\Vert \left\langle \frac{t^{2}-1}{t-1},\frac{t+1}{t}%
,t^{2}\right\rangle -\left\langle 2,2,1\right\rangle \right\Vert  &=&
\sqrt{\left( t-1\right) ^{2}\left( t^{2}+2t+3\right) } \\ 
&=&\left\vert t-1\right\vert \sqrt{\left( t+1\right) ^{2}+2} \\
&\leq& \left\vert t-1\right\vert \left( \left\vert t-1\right\vert +\sqrt{2}
\right)  \\ &<& \delta \left( \delta +\sqrt{2}\right) 
\end{eqnarray*}

si $\delta \left( \delta +\sqrt{2}\right) =\varepsilon $. Entonces 
$\delta =\frac{\sqrt{2}}{2}\left( \sqrt{2\varepsilon +1}-1\right) ;$ lo cual hace que
la definici\'{o}n se satisfaga.
\end{itemize}
Una consecuencia de la definición (\ref{preliminares5}) es la siguiente.
\begin{teorema}\label{preliminares3}
Para  $\ve{F}: A \subseteq\mathbb{R}^n \to  \mathbb{R}^m,$ 
$$ \lim_{{\bf x}\to {\bf a}}\ve{F}({\bf x})=\ve{L}=\langle l_{1},l_{2},\cdots ,l_{m}\rangle $$ existe si y solo si 
cada uno de los límites de la funciones coordenadas existe.
\end{teorema}


\begin{proof} 
La función $\ve{F}$ tiene la forma
$$\ve{F}\left( x_{1},x_{2},\cdots ,x_{n}\right) =\left\langle f_{1}\left( x_{1},x_{2},\cdots ,x_{n}\right) ,\cdots ,f_{m}\left( x_{1},x_{2},\cdots ,x_{n}\right)\right\rangle $$
\end{proof}
\begin{itemize}[leftmargin=0cm,itemindent=.2cm,labelwidth=\itemindent,labelsep=0cm,align=left]
	
\item $(\to)$ Supongamos que $\lim_{{\bf x}\to {\bf a}}\ve{F}({\bf x})= \ve{L}$, 
esto significa que para todo $\varepsilon >0$ existe $\delta >0$  tal que si 
$\left\Vert\left( x_{1},x_{2},\cdots ,x_{n}\right) -\left( a_{1},a_{2},\cdots ,a_{n}\right) \right\Vert <\delta $
entonces $\left\Vert \ve{F}\left( x_{1},x_{2},\cdots ,x_{n}\right) -\ve{L}\right\Vert <\varepsilon .$ Por lo tanto:


\begin{eqnarray*}
\left\vert f_{i}\left( x_{1},x_{2},\cdots ,x_{n}\right) -l_{i}\right\vert &=&
\sqrt{\left( f_{i}\left( x_{1},x_{2},\cdots ,x_{n}\right) -l_{i}\right) ^{2}} \\
&\leq& \sqrt{\sum_{i=1}^{m}\left( f_{i}\left(x_{1},x_{2},\cdots ,x_{n}\right) -l_{i}\right) ^{2}} \\
&=& \left\Vert \ve{F}\left( x_{1},x_{2},\cdots ,x_{n}\right) -\ve{L}\right\Vert <\varepsilon
\end{eqnarray*}
lo cual implica que $\lim_{{\bf x}\to {\bf a}} f_{i}(x_{1},x_{2},\cdots ,x_{n})=l_i, $
para $i=1,2, \cdots, m. $
\vskip0.2cm

\item $(\leftarrow)$ Ahora supongamos que $\lim_{{\bf x}\to {\bf a}} f_{i}(x_{1},x_{2},\cdots ,x_{n})=l_i,$ 
esto significa que para todo $\frac{\varepsilon}{m}>0$ existe $\delta>0$ tal que
$|f_i( x_{1},x_{2},\cdots ,x_{n}) -l_i| <\varepsilon/m ,$ siempre que
$\left\Vert\left( x_{1},x_{2},\cdots ,x_{n}\right) -\left( a_{1},a_{2},\cdots ,a_{n}\right) \right\Vert <\delta .$
Por lo tanto:
{\small
\begin{eqnarray*}
\left\Vert \ve{F}\left( x_{1},x_{2},\cdots ,x_{n}\right) - \ve{L}\right\Vert &=& 
\sqrt{\sum_{i=1}^{m}\left( f_{i}\left(x_{1},x_{2},\cdots ,x_{n}\right) -l_{i}\right) ^{2}} \\ 
&\leq & \sum_{i=1}^{m}\sqrt{\left( f_{i}\left( x_{1},x_{2},\cdots ,x_{n}\right)-l_{i}\right) ^{2}}\\
&\leq& \sum_{i=1}^{m}\left\vert f_{i}\left( x_{1},x_{2},\cdots ,x_{n}\right)-l_{i} \right\vert\\ & <& \sum_{i=1}^{m}\frac{\varepsilon }{m}=\varepsilon.
\end{eqnarray*}
}

es decir $\left\Vert \ve{F}\left( x_{1},x_{2},\cdots ,x_{n}\right) -\ve{L}\right\Vert <\varepsilon ,$
lo cual implica que $\lim_{{\bf x}\to {\bf a}}\ve{F}({\bf x})$ existe y es igual a $\ve{L}.$
\end{itemize}

\vskip0.2cm
Este resultado permite hacer cálculos directos como se ilustra a continuación.
\vskip0.2cm

\begin{itemize}[leftmargin=*]
\item {\small $\ds\lim_{t\rightarrow 0}\left\langle \frac{\sin t}{t},\sqrt{t},\frac{e^{t}-1}{t}\right\rangle= \langle \lim_{t\rightarrow 0}\frac{\sin t}{t} , \lim_{t\rightarrow 0}\sqrt{1+t} ,
\lim_{t\rightarrow 0}\frac{e^{t}-1}{t}\rangle=\langle 1 ,1,1 \ra$}
\item Un límite con un campo vectorial.
{\small
\begin{multline*} \hskip-0.2cm 
\lim\limits_{(x,y) \rightarrow \left( 1,1\right)
}\left\langle \frac{x-y}{\sqrt{x}-\sqrt{y}}, \frac{x^{3}-y^{3}}{x-y}\right\rangle 
=\lim\limits_{(x,y) \rightarrow (1,1) } \left\langle \sqrt{x}+\sqrt{y},x^{2}+xy+y^{2}\right\rangle \\
=\left\langle \lim\limits_{(x,y)
\rightarrow \left( 1,1\right) }\left( \sqrt{x}+\sqrt{y}\right)
,\lim\limits_{(x,y) \rightarrow \left( 1,1\right) }\left(
x^{2}+xy+y^{2}\right) \right\rangle = \left\langle 2,3\right\rangle 
\end{multline*}
}
{\small $$
\lim_{(x,y) \rightarrow \left( 1,1\right) }\frac{x-y}{x^{2}-y^{2}}=\lim_{(x,y) \rightarrow \left( 1,1\right) }\frac{x-y}{\left( x-y\right) \left( x+y\right) }=\lim_{(x,y)\rightarrow\left(1,1\right) }\frac{1}{\left( x+y\right) }=\frac{1}{2}
$$}
\end{itemize}

La imagen del conjunto $A$ por medio de la funci\'{o}n $\ve{F}$, se denota por $\ve{F}(A)$, es decir 
$\ve{F}(A)=\{\ve{F}(\mathbf{x})/\mathbf{x}\in A\}.$
Para funciones $\ve{F},\ve{G}:A\subseteq \mathbb{R}^n\to \mathbb{R}^m$
y $h:\mathbb{R}^n \to \mathbb{R}$ usaremos las siguientes abreviaturas: 
\begin{eqnarray*}
\left( \ve{F}\pm \ve{G}\right) \left( \mathbf{x} \right) &=& \ve{F}(\mathbf{x}) \pm \ve{G}(\mathbf{x})\\
 &=& \left\langle f_{1}(\mathbf{x}) \pm g_{1}(\mathbf{x}) ,f_{2}(\mathbf{x}) \pm g_{2}\left( 
\mathbf{x}\right) ,\cdots ,f_{m}(\mathbf{x}) \pm g_{m}\left(  \mathbf{x}\right) \right\rangle 
\end{eqnarray*}

\begin{eqnarray*}
\left( \ve{F}\cdot \ve{G}\right) \left( \mathbf{x} \right) &=&\ve{F}(\mathbf{x}) \cdot \ve{G}%
(\mathbf{x}) \\ &=& f_{1}(\mathbf{x}) g_{1}\left( \mathbf{x}\right) +f_{2}(\mathbf{x}) g_{2}(\mathbf{x})
+\cdots +f_{m}(\mathbf{x}) g_{m}(\mathbf{x}) \\
&=& \sum_{k=1}^{m}f_{k}(\mathbf{x}) g_{k}(\mathbf{x}) \\[0.1cm]
\left\Vert \ve{F}(\mathbf{x}) \right\Vert &=& \left( f_{1}^{2}(\mathbf{x}) +f_{2}^{2}(\mathbf{x}) +\cdots
+f_{m}^{2}(\mathbf{x}) \right) ^{1/2}=\sqrt{\sum_{k=1}^{m}f_{k}^{2}(\mathbf{x}) }\\[0.1cm]
h\ve{F}(\mathbf{x}) &=& h(\mathbf{x}) \ve{F}(\mathbf{x}) =\left\langle h\left( \mathbf{x} \right) f_{1}(\mathbf{x}) ,h(\mathbf{x}) f_{2}\left( \mathbf{x}\right) ,\cdots ,h(\mathbf{x}) f_{m}\left( \mathbf{x}\right) \right\rangle 
\end{eqnarray*}

Al considerar esta notación se presenta el siguiente resultado.

\begin{teorema}\label{teo1}
Sean $\ve{F},\ve{G}$ y $h$ funciones que satisfacen que 
$\ve{F},\ve{G}:A\subseteq \mathbb{R}^n\to \mathbb{R}^m$ y $h:A\subseteq \mathbb{R}^n\to \mathbb{R}$ tal que
$\lim_{\mathbf{x\to x}_{0}}\ve{F}\left( \mathbf{x}%
\right) =\mathbf{a,}$ $\lim_{\mathbf{x\to x}_{0}}\ve{G}%
(\mathbf{x}) =\mathbf{b,}$ y $\lim_{\mathbf{x\to x}%
_{0}}h(\mathbf{x}) =k;$ entonces 

\begin{enumerate}[leftmargin=0cm,itemindent=0.4cm,labelwidth=\itemindent,labelsep=0cm,align=left]

\item $\ds \lim_{\mathbf{x\to x}_{0}}\left( \ve{F}\left( \mathbf{x}%
\right) \pm \ve{G}(\mathbf{x}) \right) =\mathbf{a\pm b}$

\item $\ds\lim_{\mathbf{x\to x}_{0}}\left( h(\mathbf{x})
\ve{F}(\mathbf{x}) \right) =k\mathbf{a}$

\item $\ds\lim_{\mathbf{x\to x}_{0}}\left( \ve{F}(\mathbf{x}) \cdot \ve{G}(\mathbf{x}) \right) =\mathbf{a\cdot b}$

\item $\ds\lim_{\mathbf{x\to x}_{0}}\left\Vert \ve{F}(\mathbf{x}) \right\Vert =\left\Vert \mathbf{a}\right\Vert $
\end{enumerate}
\end{teorema}

\begin{proof}
Con la definición obtenemos las siguientes afirmaciones
\begin{enumerate}[leftmargin=0cm,itemindent=0.4cm,labelwidth=\itemindent,labelsep=0cm,align=left]
\item $\lim_{\mathbf{x\to x}_{0}}\ve{F}(\mathbf{x}) =\mathbf{a,}$
significa que para todo $\varepsilon >0$ existe $\delta >0,$ 
es decir, si $\left\Vert \mathbf{x-x}_{0}\right\Vert <\delta $ entonces 
$\left\Vert \ve{F}(\mathbf{x}) \mathbf{-a}\right\Vert <\frac{\varepsilon }{2}.$ 
Por la misma raz\'{o}n 
$\left\Vert \ve{G}(\mathbf{x}) \mathbf{-b}\right\Vert <\frac{\varepsilon }{2},$
entonces
\begin{multline*}
\left\Vert \left( \ve{F}(\mathbf{x}) \pm \ve{G}(\mathbf{x}) \right) -\left( \mathbf{a\pm b}
\right) \right\Vert = \left\Vert \left( \ve{F}\left( \mathbf{x} \right) -\mathbf{a}\right) \pm \left( \ve{G}\left( \mathbf{x} \right) -\mathbf{b}\right) \right\Vert \\
\leq \left\Vert \ve{F} (\mathbf{x}) -\mathbf{a}\right\Vert +\left\Vert \ve{G}(\mathbf{x}) -\mathbf{b}\right\Vert <\frac{\varepsilon }{2}+\frac{\varepsilon }{2}= \varepsilon 
\end{multline*}

\item Es consecuencia de la siguiente igualdad
$$h(\mathbf{x}) \ve{F}(\mathbf{x}) -k \mathbf{a=}\left( h(\mathbf{x}) -k\right) \left( \ve{%
F}(\mathbf{x}) -\mathbf{a}\right) +\left( h\left( \mathbf{x} \right) -k\right) \mathbf{a}+k\left( \ve{F}\left( \mathbf{x} \right) -\mathbf{a}\right),$$
por lo tanto,
\begin{multline*}
\left\Vert h(\mathbf{x}) \ve{F}\left( \mathbf{x}\right) -k\mathbf{a}\right\Vert \mathbf{\leq }\left\vert h\left( \mathbf{x} \right) -k\right\vert \left\Vert \ve{F}(\mathbf{x}) - \mathbf{a}\right\Vert \\ +\left\vert h(\mathbf{x}) -k\right\vert \left\Vert \mathbf{a}\right\Vert +\left\vert k\right\vert \left\Vert \ve{F}(\mathbf{x}) -\mathbf{a}\right\Vert 
\end{multline*}
Los t\'{e}rminos de la derecha tienden a cero cuando $\mathbf{x}\to \mathbf{x}_{0}.$

\item Utilizando la desigualdad de Cauchy-Schwarz, 
$0\leq \left\vert \ve{u}\cdot \ve{v}\right\vert \leq \left\Vert 
\ve{u}\right\Vert \cdot \left\Vert \ve{v}\right\Vert $
y la siguiente identidad
\begin{multline*}
\ve{F}(\mathbf{x}) \cdot \ve{G}\left( \mathbf{x}\right) -\mathbf{a\cdot b=}\left( \ve{F}\left( \mathbf{x}\right) -\mathbf{a}\right) \cdot \left( \ve{G}\left( \mathbf{x} \right) -\mathbf{b}\right) \\ 
+\mathbf{a\cdot }\left( \ve{G}\left( \mathbf{x}\right) -\mathbf{b}\right) +\mathbf{b\cdot }\left( \ve{F}(\mathbf{x}) -\mathbf{a}\right), 
\end{multline*}
se sigue que 
\begin{multline*}
\left\vert \ve{F}(\mathbf{x}) \cdot \ve{G}(\mathbf{x}) -\mathbf{a\cdot b}\right\vert \mathbf{\leq }%
\left\Vert \ve{F}(\mathbf{x}) -\mathbf{a}\right\Vert \left\Vert \ve{G}(\mathbf{x}) -\mathbf{b}\right\Vert \\
+\left\Vert \mathbf{a}\right\Vert \left\Vert \ve{G}\left( \mathbf{x}\right) -\mathbf{b}\right\Vert +\left\Vert \mathbf{b}\right\Vert \left\Vert \ve{F}(\mathbf{x}) -\mathbf{a}\right\Vert.
\end{multline*}
De las hipótesis se puede establecer que cuando $\mathbf{x\to x }_{0},$
$\left\Vert \ve{F}(\mathbf{x}) -\mathbf{a} \right\Vert \to 0,$
y $\left\Vert \ve{G}\left( \mathbf{x} \right) -\mathbf{b}\right\Vert \to 0;$ 
por lo tanto
$\left\vert \ve{F}(\mathbf{x}) \cdot \ve{G}(\mathbf{x}) -\mathbf{a\cdot b}\right\vert \to 0,$ 
cuando $\mathbf{x\to x}_{0}.$
\item El resultado es consecuencia de la siguiente desigualdad
$$\left\vert \left\Vert \ve{F}(\mathbf{x})\right\Vert -\left\Vert \mathbf{a}\right\Vert \right\vert \leq \left\Vert 
\ve{F}(\mathbf{x}) -\mathbf{a}\right\Vert $$
\end{enumerate}
\end{proof}

La siguiente es una característica especial que deseamos estudiar en la mayoría de las funciones que se consideran
en estas notas.

\subsection*{Funciones continuas}
\begin{defi} \index{Función!continua}
Se dice que una funci\'{o}n $\ve{F}:A\subseteq \mathbb{R}^n\to \mathbb{R}^m$ es continua en un punto 
$\mathbf{a}\in A$, si para todo $\varepsilon >0$ existe $\delta >0$ tal que si 
$\left\Vert \mathbf{x}-\mathbf{a}\right\Vert <\delta, $ entonces 
$\left\Vert \ve{F}(\mathbf{x}) -\ve{F}\left( \mathbf{a}\right) \right\Vert <\varepsilon .$
\vskip0.1cm
Diremos que $\ve{F}$ es continua en $A$, cuando $\ve{F}$ es continua en todo punto de $A.$
\end{defi}
Los siguientes teoremas son consecuencias de los teoremas de límites y de la definición anterior.

\begin{teorema}\label{preliminares6} Para $\ve{F}:A \subseteq\mathbb{R}^{n}\rightarrow \mathbb{R}^{m},$
$\ve{F}$ es continua si y solo si cada una sus funciones componentes es continua.
\end{teorema}

\begin{proof}
Las funciones componentes de $\ve{F}$ son campos escalares, es decir
{\small $\ve{F}\left( x_{1},x_{2},\cdots ,x_{n}\right) =\left\langle f_{1}\left(x_{1},x_{2},\cdots ,x_{n}\right) ,\cdots ,f_{m}\left( x_{1},x_{2},\cdots ,x_{n}\right)\right\rangle.$} \\
El resultado es consecuencia de aplicar el teorema (\ref{preliminares3}).
\end{proof} 

En el siguiente resultado $\ve{G}\circ \ve{F}$ denota la composición de las funciones $\ve{G}$ y $\ve{F}.$

\begin{teorema}
\normalfont Sea $\ve{F}:A\subseteq \mathbb{R}^n\to \mathbb{R}^m$ y 
$\ve{G}:B\subseteq \mathbb{R}^m\to \mathbb{R}^p$, con 
$\ve{F}\left( \mathbf{a}\right) \in B$. Si $\ve{F}$ es continua en 
$\mathbf{a}\in A$ y $\ve{G}$ es continua en $\ve{F}\left( \mathbf{a}\right) $ 
entonces $\ve{G}\circ \ve{F}:A\subseteq \mathbb{R}^n\to \mathbb{R}^p$ es continua en $\mathbf{a}$.
\end{teorema}

\begin{proof}
Por ser $\ve{G}$ continua en $\ve{F}\left( \mathbf{a}\right) \in B,$ dado $\varepsilon >0$ 
existe $\lambda >0,$ para $\mathbf{y} \in B,$ si $\left\Vert \mathbf{y-}\ve{F}\left( \mathbf{a}\right)
\right\Vert <\lambda $ entonces $\left\Vert \ve{G}(\mathbf{y})-\ve{G}\left( \ve{F}\left( \mathbf{a}\right) \right) \right\Vert <\varepsilon .$
\vskip0.1cm 
Como $\ve{F}$ continua en $\mathbf{a}\in A,$ esto implica que dado 
$\delta >0,$ existe $\lambda >0,$  si $\left\Vert \mathbf{x-a}\right\Vert <\delta ,$ para $\mathbf{x}\in A,$
entonces $\left\Vert \ve{F}(\mathbf{x}) -\ve{F}\left( \mathbf{a}\right) \right\Vert <\lambda ,$ 
lo anterior implica que 
$\left\Vert \ve{G}\left( \ve{F}\left( \mathbf{x} \right) \right) -\ve{G}\left( \ve{F}\left( 
\mathbf{a}\right) \right) \right\Vert <\varepsilon .$
En síntesis, dado $\varepsilon >0,$ existe $\delta >0$ tal que si 
$\left\Vert \mathbf{x-a}\right\Vert <\delta $ entonces 
$\left\Vert \ve{G} \left( \ve{F}(\mathbf{x}) \right) -\ve{G}\left( \ve{F}\left( \mathbf{a}\right) \right)
\right\Vert <\varepsilon ;$ es decir, $\ve{G}\circ \ve{F}$ es continua en $\mathbf{a.}$
\end{proof} 

\begin{teorema}
Sea $\ve{F},\ve{G}:A\subseteq \mathbb{R}^n\to \mathbb{R}^m$ y $h:A\subseteq \mathbb{R}^n\to \mathbb{R}$, 
funciones continuas en $\mathbf{a}\in A,$ entonces las siguientes funciones son continuas en $\mathbf{a}$ 
%
$\ve{F}\pm \ve{G}:A\subseteq \mathbb{R}^n\to \mathbb{R}^m,$ \\
$ \ve{F}\cdot \ve{G}:A\subseteq \mathbb{R}^n\to \mathbb{R},$\\
$\left\Vert  \ve{F}\right\Vert :A\subseteq \mathbb{R}^n\to \mathbb{R},$\\
$h\ve{F}:A\subseteq \mathbb{R}^n\to \mathbb{R}^m.$ 
%
\end{teorema}
La prueba es consecuencia de aplicar los teoremas (\ref{preliminares5}) y (\ref{teo1}).
En particular, el lector puede verificar la continuidad de la función ${\bf u}(t) \times {\bf v}(t)$ para dos funciones vectoriales continuas 
${\bf u}, {\bf v}: \mathbb{R}\to \mathbb{R}^3.$

\section{Diferenciabilidad} \label{diferencibilidad}
Para funciones reales $f:\mathbb{R}\rightarrow \mathbb{R}$, el lector seguramente
estar\'{a} familiarizado con la definici\'{o}n de la derivada en un punto $a$,
$$f^{\prime}(a) =\lim_{h\rightarrow 0}\frac{f(a+h)-f(a) }{h}.$$
Al considerar la transformaci\'{o}n lineal $A(h)=f^{\prime}(a)\cdot h.$
Entonces, la expresi\'{o}n anterior se reescribe como 
$$\lim_{h\rightarrow 0}\frac{f\left( a+h\right) -f\left( a\right) -A\left( h\right) }{h}=0.$$

En el caso general una funci\'{o}n ${\bf F}:\mathbb{R}^{n}\rightarrow \mathbb{R}^{m}$ se dice
diferenciable en $a\in \mathbb{R}^{n}$
si existe una transformaci\'{o}n lineal $\lambda :\mathbb{R}^{n}\rightarrow \mathbb{R}^{m},$ tal que 
$$\lim_{h\rightarrow 0}\frac{\left\Vert {\bf F}(a+h)-{\bf F}(a)-\lambda(h)\right\Vert }{\left\Vert h\right\Vert }=0.$$
Es de anotar que $\left\Vert {\bf F}\left( a+h\right) -{\bf F}(a) -\lambda(h) \right\Vert$ representan la norma en 
$\mathbb{R}^{m}$ y $\left\Vert h\right\Vert $ es la norma en $\mathbb{R}^{n}.$ \index{Función!diferenciable}
Se dice que ${\bf F}$ es diferenciable en $A$, si es diferenciable para todo $a\in A.$
La mencionada transformación lineal es la denominada \emph{matriz jacobiana} de ${\bf F}$ en $a$,
que se denota por $D({\bf F})(a)$. Para ${\bf F}:\mathbb{R}^n\to {\mathbb R}^m$ y $a=(a_1,a_2,\cdots, a_n)$, es claro que
${\bf F}(a)=\langle f_1(a),f_2(a), \cdots, f_m(a) \rangle. $
La entrada $ij$ de la matriz jacobiana está dada por
$$\frac{\partial f_i}{\partial x_j}(a)=\lim_{h\rightarrow 0}\frac{f_i(a_1,a_2,\cdots,a_j+h,\cdots,a_n)-f_i(a_1,a_2,\cdots,a_j,\cdots, a_n)}{h}.$$
que la denominaremos la derivada parcial de $f_i$ con respecto a $x_j$. \index{Derivada parcial}
La matriz jacobiana de ${\bf F}$, $D{\bf F}(a)$ es 
$$D{\bf F}(a)=\left(\begin{array}{cccc}
\frac{\partial f_1}{\partial x_1}(a)&\frac{\partial f_1}{\partial x_2}(a)&\cdots & \frac{\partial f_1}{\partial x_n}(a)\\[0.2cm]
\frac{\partial f_2}{\partial x_1}(a)&\frac{\partial f_2}{\partial x_2}(a)&\cdots & \frac{\partial f_2}{\partial x_n}(a)\\[0.2cm]
\vdots \\[0.2cm]
\frac{\partial f_m}{\partial x_1}(a)&\frac{\partial f_m}{\partial x_2}(a)&\cdots & \frac{\partial f_m}{\partial x_n}(a)
\end{array} \right)$$
\begin{teorema}
Si ${\bf F}:\mathbb{R}^n \to \mathbb{R}^m$ es diferenciable en $a$ y 
${\bf G}:\mathbb{R}^m \to \mathbb{R}^p$ es diferenciable en ${\bf F}(a)$ entonces la función compuesta
${\bf G}\circ {\bf F}:\mathbb{R}^n \to \mathbb{R}^p$ es diferenciable en $a$. Además
$D({\bf G}\circ {\bf F})(a)=D({\bf G}( {\bf F}(a))\cdot D{\bf F}(a). $
\end{teorema}
La prueba de este resultado puede consultarse en \textcite{spivak_calculus_2018}. 
En el caso de campos escalares $f:\mathbb{R}^n \to \mathbb{R}$ a $Df(a)$ se denota también por 
$\nabla f(a)$ y se denomina el \emph{vector gradiente} de $f$. \index{Gradiente}
\vskip0.2cm
Un campo vectorial ${\bf F}(x,y,z)$ con dominio $D \subseteq \mathbb{R}^3$ 
se denomina \emph{conservativo} si existe un campo escalar
$\phi(x,y,z)$ que se denomina \emph{función potencial} con dominio $D$ y satisface que $\nabla \phi ={\bf F}.$ \index{Campo vectorial!conservativo} \index{Función!potencial}
En \cite{thomas_calculo_2015} se presenta el siguiente criterio, el cual es útil para determinar si un campo vectorial cumple esta condición.
\begin{prop}
Sea ${\bf F} = P(x, y,z)\vi +Q(x,y,z)\vj +P(x,y,z)\vk$ un campo en un dominio abierto simplemente
conexo, cuyas funciones componentes tienen sus primeras derivadas parciales continuas. 
De esta forma, ${\bf F}$ es conservativo si y sólo se satisfacen la siguientes condiciones:
$$\frac{\partial P}{\partial y}=\frac{\partial Q}{\partial x},\hspace{1.5cm}
\frac{\partial P}{\partial z}=\frac{\partial R}{\partial x},\hspace{1.5cm}
\frac{\partial Q}{\partial z}=\frac{\partial R}{\partial y}.$$
\end{prop}
\section{Funciones vectoriales} \index{Funciones vectoriales}
De manera general, una aplicación que asigna a cada número real un vector en $\mathbb{R}^n$
se denomina función vectorial, esto es ${\bf r}(t)= \langle r_1(t),r_2(t),\cdots, r_n(t)\ra$.
\vskip0.2cm
Por ejemplo, la función ${\bf r}(t)= \langle 3\cos(2t),1-\cos(t),2\sin(t)\ra$ está definida para $t\in [0,2\pi]$ su 
gráfico en $\mathbb{R}^3$ se representa a continuación.


\vskip0.2cm
\noindent
\begin{minipage}{\textwidth}
	\centering
	\captionof{figure}{Curva determinada por ${\bf r}(t)= \langle 3\cos(2t),1-\cos(t),2\sin(t)\ra$}
	\begin{minipage}{0.38\linewidth}
		\centering
		\includegraphics[width=\linewidth]{Fig/4.pdf}
	\end{minipage}
	\begin{minipage}{0.38\linewidth}
		\centering
		\includegraphics[width=\linewidth]{Fig/5.pdf}
	\end{minipage}
	
	\vspace{0.3em}
	\parbox{0.95\linewidth}{\centering\fontsize{8pt}{10pt}\selectfont\textbf{Fuente:} elaboración propia.}
\end{minipage}
\vskip0.2cm


Para una función vectorial ${\bf r}(t)=\langle r_1(t),r_2(t),\cdots, r_m(t) \rangle ,$
los dos primeros de los siguientes resultados se obtienen 
como caso particular de los teoremas (\ref{preliminares3}) y (\ref{preliminares6}), con $n=1.$ 

\begin{itemize}
\item $\lim_{t\to a}{\bf r}(t)$ existe si y solo si $\lim_{t\to a}r_i(t)$ para toda $i=1,2,\cdots ,m$.
\item ${\bf r}(t)$ es continua en $a$ si y solo si cada una de las funciones componentes $r_i(t)$ es continua en $a$, para toda $i=1,2,\cdots ,m$.
\item ${\bf r}(t)$ es diferenciable en $a$ si y solo si cada una de las funciones componentes $r_i(t)$ es diferenciable en $a$, para toda $i=1,2,\cdots ,m$; además ${\bf r}^\prime(t)= \langle r^\prime_1(t),r^\prime_2(t),\cdots, r^\prime_m(t)\ra.$
\end{itemize}

Presentamos la demostración del último de estos resultados.
\begin{eqnarray*}
\mathbf{r}^{\prime}(t) &=& \lim_{h\rightarrow 0}\frac{1}{h}\left( \mathbf{r}\left( t+h\right) -\mathbf{r}(t) \right) \\
&=& \lim_{h\rightarrow 0}\frac{1}{h}\left\langle r_{1}\left(t+h\right) -r_{1}(t) ,\cdots ,r_{m}\left( t+h\right) -r_{m}(t) \right\rangle \\
&=&\left\langle \lim_{h\rightarrow 0}\frac{r_{1}\left( t+h\right) -r_{1}(t) }{h},\cdots ,\lim_{h\rightarrow 0}\frac{r_{m}\left( t+h\right) -r_{m}(t) }{h}\right\rangle \\
&=& \left\langle r_{1}^{\prime}(t),r_{2}^{\prime}(t) ,\cdots ,r_{m}^{\prime}\left(t\right) \right\rangle 
\end{eqnarray*}

Si $\mathbf{u}(t) =\left\langle u_{1}(t),u_{2}(t) ,\cdots ,u_{n}(t) \right\rangle $ y 
$\mathbf{v}(t) =\left\langle v_{1}(t) ,v_{2}\left(t\right) ,\cdots ,v_{n}(t) \right\rangle $
son dos funciones vectoriales entonces el producto punto entre funciones vectoriales está dado por
$\mathbf{u}(t) \cdot \mathbf{v}(t) =u_{1}\left(t\right) v_{1}(t) +u_{2}(t) v_{2}(t)+\cdots +u_{n}(t) v_{n}(t) .$
Derivando se obtiene las siguientes igualdades
\begin{eqnarray*}
\left( \mathbf{u}(t) \cdot \mathbf{v}(t) \right)^{\prime} 
%&=& u_{1}^{\prime(t) v_{1}(t) +u_{1}(t) v_{1}^{\prime }(t) +u_{2}^{\prime }(t)v_{2}(t) +u_{2}(t) v_{2}^{\prime }(t)+ \cdots + u_{n}^{\prime }(t) v_{n}(t) +u_{n}\left(t\right) v_{n}^{\prime }(t) \\
&=& \big( u_{1}^{\prime }(t) v_{1}(t) +u_{2}^{\prime}(t) v_{2}(t) +\cdots +u_{n}^{\prime }(t) v_{n}(t) \big) + \\ 
&& \big( u_{1}(t)v_{1}^{\prime }(t) +  u_{2}(t) v_{2}^{\prime }(t) +\cdots +u_{n}(t) v_{n}^{\prime }(t)\big) \\
&=& \mathbf{u}^{\prime }(t) \cdot \mathbf{v}(t) +\mathbf{u}(t) \cdot \mathbf{v}^{\prime }(t) 
\end{eqnarray*}
Lo cual se resume en la siguiente fórmula
\begin{equation}
\left( \mathbf{u}(t) \cdot \mathbf{v}(t) \right)^{\prime }=\mathbf{u}^{\prime }(t) \cdot \mathbf{v}\left(
t\right) +\mathbf{u}(t) \cdot \mathbf{v}^{\prime }\left(t\right).
\end{equation}
Otras fórmulas útiles, son las siguientes
\begin{eqnarray*}
(\left\Vert \mathbf{u}(t) \right\Vert)^{\prime } &=& \frac{d}{dt}\left( \left( \mathbf{u}(t) \cdot \mathbf{u}(t)
\right) ^{1/2}\right) \\
&=& \frac{1}{2}\left( \mathbf{u}(t) \cdot \mathbf{u}(t) \right) ^{-1/2}\left( \mathbf{u}^{\prime }(t) \cdot 
\mathbf{u}(t) +\mathbf{u}(t) \cdot \mathbf{u}^{\prime }(t) \right) = \frac{\mathbf{u}^{\prime }(t)\cdot \mathbf{u}(t) }{\left\Vert \mathbf{u}(t)\right\Vert }
\end{eqnarray*}

La derivada del producto cruz entre dos funciones vectoriales se evalua como sigue
\begin{prop}
Si $\mathbf{u}(t)$ y $\mathbf{v}(t)$ son funciones vectoriales definidas en $\mathbb{R}^3$ entonces
$\left(\mathbf{u}(t)\times\mathbf{v}(t)\right)^{\prime}= \mathbf{u}^{\prime}(t)\times\mathbf{v}(t)+\mathbf{u}(t)\times \mathbf{v}^{\prime}(t)$.
\end{prop}

\begin{proof}
Si $\mathbf{u}(t)=\left\langle u_{1}(t),u_{2}(t),u_{3}(t) \right\rangle $ y 
$\mathbf{v}(t)=\left\langle v_{1}(t),v_{2}(t),v_{3}(t)\right\rangle ,$
obviaremos el uso de la variable $t$ para hacer menos extensa la escritura, es decir
\begin{eqnarray*}
\frac{d}{dt}\left( \mathbf{u}(t)\times \mathbf{v}(t)\right) &=& \frac{d}{dt} \langle u_{2} v_{3} -u_{3} v_{2} ,u_{3} v_{1} -u_{1} v_{3} ,u_{1} v_{2} -u_{2} v_{1} \rangle \\
&=& \big\langle u_{2}^{\prime} v_{3}+u_{2} v_{3}^{\prime} -u_{3}^{\prime} v_{2} -u_{3} v_{2}^{\prime},u_{3}^{\prime} v_{1} +u_{3}v_{1}^{\prime} \\
&& -u_{1}^{\prime}v_{3} -u_{1} v_{3}^{\prime},   u_{1}^{\prime}v_{2}+u_{1}v_{2}^{\prime}  -u_{2}^{\prime} v_{1} -u_{2} v_{1}^{\prime} \big\rangle \\
&=& \left\langle u_{2}^{\prime} v_{3}-u_{3}^{\prime} v_{2} ,u_{3}^{\prime}(t) v_{1} -u_{1}^{\prime} v_{3},u_{1}^{\prime} v_{2} -u_{2}^{\prime} v_{1} \right\rangle  \\
&& +\left\langle u_{2} v_{3}^{\prime} -u_{3}v_{2}^{\prime} ,u_{3} v_{1}^{\prime}(t) -u_{1} v_{3}^{\prime} ,u_{1} v_{2}^{\prime} -u_{2} v_{1}^{\prime} \right\rangle \\
&=& \mathbf{u}^{\prime}(t)\times \mathbf{v}(t)+\mathbf{u}(t)\times \mathbf{v}^{\prime}(t) 
\end{eqnarray*}
\end{proof}
La siguiente expresión es producto de aplicar los resultados anteriores
\begin{multline}\left( \left( \mathbf{u}(t)\times \mathbf{v}(t)\right)\cdot \mathbf{w}(t)\right)^{\prime} \\ = \left( \left( \mathbf{u}^{\prime}(t)\times \mathbf{v}(t)+\mathbf{u}(t)\times \mathbf{v}^{\prime}(t)\right) \cdot \mathbf{w}(t)\right) +\left( \left( \mathbf{u}(t)\times \mathbf{v}(t)\right)\cdot \mathbf{w}^{\prime}(t)\right) 
\end{multline}

\section{Regiones parametrizadas} \label{regionesparametrizadas}
En estas notas trabajaremos básicamente sobre tres tipos de objetos en el espacio: curvas, superficies y sólidos.
En esta sección presentaremos una manera de representar dichos objetos, la idea consiste en \emph{parametrizar} la región en mención. Para este propósito se utiliza, según el caso, uno, dos o tres parámetros, respectivamente.  
\index{Parametrizar}

\vspace{-1mm}

\begin{tcolorbox}[breakable,enhanced,title=\bf Curva parametrizada]
La funci\'on $\ve{r}(t)$ se denomina una \emph{parametrizaci\'on} de una curva $C$ si 
al variar el parametro $t\in [a,b]$ la imagen $\ve{r}(t)$ describe los puntos de la curva $C$. 
Si $\ve{r}(a)=\ve{r}(b)$ la curva se denomina cerrada y la curva se denomina simple si $\ve{r}(t)\neq \ve{r}(s),$
para $a\neq t<s\neq b$.
\end{tcolorbox}

Las funciones vectoriales son un ejemplo de parametrización de curvas.
El concepto de curva simple hace referencia a que la curva no se autointerseca, posiblemente solo en los puntos extremos.
La curva se denomina \emph{suave} si $\ve{r}^\prime(t) \neq \ve{0},$ para todo $t.$ \index{Curva!suave}

\begin{tcolorbox}[breakable,enhanced,title=\bf Superficie parametrizada]
Una \emph{parametrización} de una superficie $\sigma $
es una aplicación $$\ve{r}:[a,b]\times \lbrack c,d]\rightarrow \mathbb{R}^{3}$$
tal que al variar los parámetros $u\in \lbrack a,b]$ y $v\in \lbrack c,d]$
su imagen $\ve{r}(u,v)$ va describiendo los puntos de la superficie $\sigma $.
\end{tcolorbox}
La superficie se denomina \emph{suave} si $\frac{\partial \ve{r}}{\partial u}$ y 
$\frac{\partial \ve{r}}{\partial v} $ son continuas y 
$\frac{\partial \ve{r}}{\partial u} \times \frac{\partial \ve{r}}{\partial u} \neq \ve{0}.$ 
\index{Superficie!suave}
%\vskip0.2cm
Con los elementos anteriormente descritos es posible construir funciones
que describan un sólido en el espacio tridimensional. 
\begin{tcolorbox}[breakable,enhanced,title=\bf Parametrización de un sólido]
Se dice que la función $${\bf r}(u,v,w)= \langle x(u,v,w), y(u,v,w),z(u,v,w) \rangle $$  
parametriza el sólido $E$ si al variar los parámetros se recorre completamente a $E$,  
con $a \leq u \leq b$, $c \leq v \leq d$ y $p \leq u \leq q$.  
\end{tcolorbox}


Una ventaja de usar este tipo de representación es el hecho de que al usar límites constantes el planteamiento de integrales
resulta, en algunos casos, bastante sencillo.

\section{Elementos básicos de integración}

Presentamos en esta sección algunos de los elementos teóricos propios de un curso de 
Cálculo Vectorial que serán usados en los capítulos posteriores. La exposición teórica completa puede ser consultada 
por el lector con más detalle en \cite{apostol_alisis_1981}, \cite{marsden_calculo_2013}, \cite{stewart_calculo_2018} o \cite{thomas_calculo_2015}.
\vskip0.2cm
Posteriormente se utilizan estos elementos en algunas aplicaciones de la integración principalmente en 
el c\'{a}lculo de longitud de curvas, \'{a}reas, volumenes y centros de masa. 
Se estudiar\'{a} tambi\'{e}n el teorema de cambio de variables,
particularmente cuando se utilizan coordenadas polares, cilindricas o esf\'{e}ricas, pero se har\'{a} enfasis en 
parametrizar la regi\'{o}n de integraci\'{o}n.
\vskip0.2cm
Una función $g:\mathbb{R}\to\mathbb{R}$ se dice \emph{integrable} en $[a,b]$ en el sentido señalado en \textcite{apostol_calculus_1999};
utilizaremos esta idea en los siguientes resultados. 

\subsection*{Integrales de línea en campos escalares}

Consideremos una curva $\alpha$ parametrizada por una función 
${\bf r}:\mathbb{R}\to\mathbb{R}^3$ continua, inyectiva y derivable en $(a,b)$, con ${\bf r}(t)=\langle x(t),y(t),z(t) \rangle .$ 
% \index{Partición}
\vskip0.2cm

\noindent
\begin{minipage}{\textwidth}
	\centering
	\captionof{figure}{Curva en $\mathbb{R}^3$}
	\label{Fig/6.png}
	\begin{minipage}{0.32\linewidth}
		\centering
		\includegraphics[width=\linewidth]{Fig/6.png}
	\end{minipage}
	
	\vspace{0.3em}
	\parbox{0.95\linewidth}{\centering\fontsize{8pt}{10pt}\selectfont\textbf{Fuente:} elaboración propia.}
\end{minipage}

\vskip0.2cm
Una \emph{partición} $P$ del intervalo $[a,b]$ es un conjunto finito de puntos 
$P=\{ t_0, t_1, t_2, \cdots , t_{n-1},t_n \}$ que satisface 
$$a = t_0< t_1<  t_2< \cdots < t_{n-1}< t_n=b;$$
la cual determina una partición de la curva $\alpha$ en $n$ puntos de la forma
$Q_i(x(t_{i-1}), x(y_{i-1}),z(t_{i-1})),$ para $i=1,2, \cdots , n.$
%

Denotaremos por $$\Delta\alpha_i=\int_{t_i}^{t_{i+1}}\Vert {\bf r}^\prime(t)\Vert dt $$ 
a la longitud de la curva $\alpha$ desde el punto $Q_i$ hasta $Q_{i+1}$.
Para una función continua $f: \alpha \to \mathbb{R}$ consideremos la suma 
$$\sum_{i=1}^n f({\bf r}(t_i^*))\Delta \alpha_i=\sum_{i=1}^n f(x(t_i^*),y(t_i^*),z(t_i^*))\Delta \alpha_i$$
en donde $t_i^* \in [t_{i-1},t_i],$ con los elementos anteriores se define la integral de una campo escalar sobre una curva.
\begin{defi}
Consideremos $f:\mathbb{R}^3\to \mathbb{R}$ continua. 
La integral de $f$ sobre la curva suave $\alpha$ con respecto a la parametrización 
${\bf r}(t)$ está dada por
\begin{equation} \label{intdelinea}
\int_a^b f({\bf r}(t))\Vert{\bf r}^\prime(t)\Vert dt=\lim_{n\to\infty}\left(\sum_{i=1}^n f({\bf r}(t_i^*))\Delta \alpha_i\right) \end{equation} \index{Integrales!de línea}
\end{defi}
En particular, si la función $f$ de la integral (\ref{intdelinea}) es $1,$ obtendremos la fórmula para la longitud de la 
curva (35) y si $f$ es la densidad $\delta(x,y,z)$ de un alambre o varilla delgada cuya forma está determinada por la curva $\alpha$ en el punto $(x,y,z),$ entonces la masa $M$ de la varilla está dada por
\begin{equation} 
M=\int_a^b \delta({\bf r}(t))\Vert{\bf r}^\prime(t)\Vert dt \end{equation}
Los momentos con respecto a los planos coordenados se definen como
\begin{eqnarray*} 
\mu_{xy}&=&\int_a^b z(t)\,\delta({\bf r}(t))\Vert{\bf r}^\prime(t)\Vert dt,\\
\mu_{yz}&=&\int_a^b x(t)\,\delta({\bf r}(t))\Vert{\bf r}^\prime(t)\Vert dt,\\
\mu_{xz}&=&\int_a^b y(t)\,\delta({\bf r}(t))\Vert{\bf r}^\prime(t)\Vert dt.
\end{eqnarray*}
Las coordenadas del centro de masa de la varilla son 
\index{Centro de masa} 
$$(\overline{x},\overline{y},\overline{z})= (\frac{\mu_{yz}}{M},\frac{\mu_{xz}}{M},\frac{\mu_{xy}}{M}).$$

\subsection*{Integración de funciones vectoriales} \index{Función!integrable}
Para funciones vectoriales ${\bf f}(t)=\langle f_1(t),f_2(t),\cdots,f_n(t)\ra,$ se dice que $f$ es integrable en $[a,b]$ 
si y solo si cada una de las funciones componentes $f_i(t)$ es integrable en $[a,b]$, para toda $i=1,2,\cdots ,n$.  
Además $$\int_a^b{\bf f}(t)dt= \langle \int_a^b f_1(t)dt,\int_a^b f_2(t)dt,\cdots, \int_a^b f_n(t)dt\ra.$$

\subsection*{Integrales dobles} \label{integralesdobles}  \index{Integrales!dobles}
Presentaremos ahora la integral doble de una funci\'{o}n de
dos variables $f(x,y)$ definida sobre una regi\'{o}n del plano como el 
l\'{\i}mite de las aproximaciones a trav\'{e}s de sumas de Riemann, posteriormente
se presentar\'{a}n de manera similar las integrales triples. 
La evaluaci\'{o}n de las integrales dobles o triples se puede obtener
utilizando el teorema fundamental de c\'{a}lculo; véase \cite{apostol_calculus_1999}.
\vskip0.2cm
En el plano $xy$ consideraremos un rectángulo $R=[a,b]\times \lbrack c,d],$ es decir
$$R=\{(x,y) \in \mathbb{R}^{2}\, | \,a\leq x\leq b,c\leq y\leq d\}.$$
Si $f(x,y)$ es una funci\'{o}n continua definida sobre el rect\'{a}ngulo $R,$ 
subdividiremos a $R$ en $nm$ rect\'{a}ngulos m\'{a}s peque\~{n}os a partir
de rectas paralelas a los ejes coordenados. Esto es, para el intervalo $[a,b],$ consideremos
$$a=x_{0}<x_{1}<x_{2}<\cdots <x_{i}<\cdots <x_{n-1}<x_{n}=b,$$ mientras que
para el intervalo $[c,d],$ consideremos
$$c=y_{0}<y_{1}<y_{2}<\cdots <y_{k}<\cdots <y_{m-1}<y_{m}=d.$$ A esta división de $R$, lo denominaremos una partición de $R$.
\index{Partición}
Si $n$ y $m$ crece entonces el ancho y el largo de estos subintervalos decrecen.
\vskip0.2cm
Denotaremos $\Delta x_{i}$ al ancho del $i-$\'{e}simo subintervalo
sobre el eje $x,$ y $\Delta y_{k}$ es el alto $k-$\'{e}simo subintervalo sobre el eje $y$, 
es decir $\Delta x_{i}=x_{i+1}-x_{i},$ de la misma forma $\Delta
y_{k}=y_{k+1}-y_{k}.$ El \'{a}rea de uno de estos subintervalos de ancho $\Delta x_{i}$ y altura $\Delta y_{k}$ es
$\Delta A_{ik}=\Delta x_{i}\Delta y_{k}$.
La norma de la partici\'{o}n se denota por $\left\Vert P\right\Vert$ y es el m\'{a}ximo sobre todos los valores 
$i$ y $k$, de todas las dimensiones de los rect\'{a}ngulos en que se divide la regi\'{o}n $R.$
\vskip0.2cm
Cuando todos estos subintervalos son de igual longitud, los denotaremos por $\Delta x=\frac{b-a}{n}$ y $\Delta y=\frac{d-c}{m}.$
\vskip0.2cm \index{Suma doble de Riemann}
Una suma de Riemann de $f(x,y) $ sobre el rect\'{a}ngulo $R$ se obtiene al considerar un punto cualquiera 
$\left( x_{i}^{\ast },y_{k}^{\ast}\right) ,$ en $ik$-\'{e}simo subrect\'{a}ngulo, esto 
$\left( x_{i}^{\ast},y_{k}^{\ast }\right) \in \lbrack x_{i-1},x_{i}]\times \lbrack y_{k-1},y_{k}],$ 
o de manera similar $x_{i-1}<x_{i}^{\ast }<x_{i},$ y $y_{k-1}<y_{k}^{\ast }<y_{k};$  
El valor de la funci\'{o}n es este punto es $f\left( x_{i}^{\ast},y_{k}^{\ast }\right) $ y este valor se 
multiplica por el \'{a}rea del $ik$-\'{e}simo rect\'{a}ngulos; se suman todos estos productos para obtener una
expresi\'{o}n de la forma
$$S_{nm}=\sum_{i=1}^{n}\sum_{k=1}^{m}f\left( x_{i}^{\ast },y_{k}^{\ast}\right) \Delta x_{i}\Delta y_{k}.$$

Dependiendo de la selecci\'{o}n $\left( x_{i}^{\ast },y_{k}^{\ast }\right) ,$
es posible obtener diversos valores de $S_{nm}.$
\vskip0.2cm
Es de especial inter\'{e}s el conocer el comportamiento de las sumas de
Riemann cuando la norma de la partici\'{o}n tiende a cero, que equivale a
que el ancho y la altura de cada subrect\'{a}ngulo de $R$ tienden a cero.
Si el l\'{\i}mite de las sumas $S_{nm}$ existe sin importar la selecci\'{o}n
que se haga de $\left( x_{i}^{\ast },y_{k}^{\ast }\right) ,$ entonces
diremos que la funci\'{o}n $f(x,y) $ es integrable sobre el rect\'{a}ngulo $R$ 
y este valor se denota por 
$\ds \iint\limits_{R}f\left(x,y\right) dA$ o tambi\'{e}n $\ds \iint\limits_{R}f(x,y) dydx.$
Lo anterior nos motiva a precisar este concepto.

\begin{defi}
La integral doble de una función continua $f(x,y)$ sobre el rectángulo $R$ se define como
\begin{eqnarray} \label{integralesdobles1}
\lim _{\left\Vert P\right\Vert \to 0} S_{nm} &=& \lim _{n,m \to \infty}\left(\sum_{i=1}^{n}\sum_{k=1}^{m}f\left( x_{i}^{\ast },y_{k}^{\ast}\right) \Delta x_{i}\Delta y_{k}\right) \nonumber \\ &=& \ds \iint\limits_{R}f(x,y) dA
\end{eqnarray}
cuando este límite existe.
\end{defi}

\subsection*{Aplicaciones de las integrales dobles}

Si la funci\'{o}n $f(x,y)$ es positiva sobre el rect\'{a}ngulo $R$ del plano $xy$ la integral doble de $f$ sobre $R$ puede
interpretarse como el volumen del s\'{o}lido $E$ acotado en la base por $R,$
por la parte superior la superficie $z=f(x,y),$ y en los lados regiones contenidas en
los planos $x=a,$ $x=b$, $y=c$, $y=d.$ N\'{o}tese que el t\'{e}rmino
$f(x_{i}^{\ast },y_{k}^{\ast })\Delta x_{i}\Delta y_{k}$ corresponde al
volumen del paralelep\'{\i}pedo con \'{a}rea de la base $\Delta x_{i}\Delta y_{k}$ y altura $f(x_{i}^{\ast },y_{k}^{\ast });$
la suma $S_{nm}$ es una aproximaci\'{o}n al volumen del s\'{o}lido $E.$
Es por lo anterior que el volumen de $E$ se define como 
$$\text{volumen:}\lim _{n,m \to \infty}\left(\sum_{i=1}^{n}\sum_{k=1}^{m}f\left( x_{i}^{\ast },y_{k}^{\ast}\right) \Delta x_{i}\Delta y_{k}\right)= \ds \iint\limits_{R}f(x,y) dA.$$ \index{Masa de una región}

Si la función $f(x,y)$ determina la densidad en el punto $(x,y)$ entonces la integral de la fórmula (\ref{integralesdobles1})
representa la masa ($M$) de la región $R$. Los momentos con respecto al eje $x$ y al eje $y$ están dados  por
{\small $\mu_{x}=\ds \iint_R y f(x,y)dA$} y {\small $\mu_{y}=\ds \iint_R x f(x,y)dA$,} los cuales determinan las coordenadas del 
centro de masa de la región $R$, esto es $$(\overline{x},\overline{y})= \left(\frac{\mu_y}{M},\frac{\mu_x}{M}\right).$$
\index{Centro de masa}

Si la región de integración en la fórmula (\ref{integralesdobles1}) es un rectángulo está se evalúa
$$\ds \iint\limits_{R}f(x,y) dA= \int_c^d \left (\int_a^b f(x,y) dx \right )dy.$$
\vskip0.2cm
El teorema de Fubini, cuya demostración puede consultarse en \cite{marsden_calculo_2013} o \cite{apostol_alisis_1981},
es el resultado que permite intercambiar el orden de integración incluso en regiones más generales. 
\vskip0.2cm
Si la región de integración no es un rectángulo, 
la integral doble de una función $f(x,y)$ sobre una región $R$ puede evaluarse mediante la aplicación del Teorema de 
Fubini por alguna de las siguientes expresiones, dependiendo de la forma de la región $R$.
\vskip0.2cm
Si la región del plano en mención, se describe en la forma $\phi_{1}(x) \leq y \leq \phi_{2}(x)$ con $a \leq x \leq b$, 
diremos que la región es de tipo I; mientras que la región del tipo II es aquella que puede ser descrita como
$\varphi_{1}(y) \leq x \leq \varphi_{2}(y)$ con $c \leq y \leq d.$ 

\vskip0.2cm

De acuerdo con lo anterior, el planteamiento de la integral doble es, según el caso, el siguiente:

\vskip0.2cm
\noindent
\begin{minipage}{\textwidth}
	\centering
	\captionof{figure}{Regiones de integración en el plano}
	\label{Regiones_integracion_plano}
	\begin{minipage}{0.35\linewidth}
		\centering
		\includegraphics[width=\linewidth]{Fig/7.pdf}
	\end{minipage}
	\begin{minipage}{0.35\linewidth}
		\centering
		\includegraphics[width=\linewidth]{Fig/8.pdf}
	\end{minipage}
	
	\vspace{0.3em}
	\parbox{0.95\linewidth}{\centering\fontsize{8pt}{10pt}\selectfont\textbf{Fuente:} elaboración propia.}
\end{minipage}
\vskip0.2cm



$$ \iint_{R}f(x,y) dA=\int_{a}^{b}\int_{\phi _{1}(y) }^{\phi_{2}(y)}f(x,y)dydx,$$

o también
\vskip0.2cm
$$ \iint_{R}f(x,y)dA=\int_{c}^{d}\int_{\varphi _{1}(y) }^{\varphi _{2}(y)}f(x,y)dxdy.$$

%\vskip0.2cm

\begin{ejemplo}
Consideremos la función $f(x,y)=8+2x-x^3-4y$ y algunas aproximaciones para diferentes valores de $n$ y $m$;
con $x\in [0,2]$, $y \in [0,1].$

%\captionsetup{hypcap=false}
%\captionof{figure}{\small Aproximaciones al volumen bajo $f(x,y)$}
\vspace{4mm}
\noindent % Evita sangrías no deseadas
\begin{minipage}{\textwidth}
	\centering
	% Usamos \captionof porque estamos fuera de un entorno figure
	\captionof{figure}{Aproximaciones para el volumen bajo $f(x,y) = 8+2x-x^3-4y$.}
	\label{fig:aproximaciones_volumen}
	
	\vspace{0.3cm} % Espacio entre título e imágenes
	
	\begin{tabular}{ccc}
		% PRIMERA FILA: Las imágenes con escalas homogeneizadas
		\includegraphics[scale=0.31]{Img/T_51a.pdf} & 
		\includegraphics[scale=0.28]{Img/T_51b.pdf} &
		\includegraphics[scale=0.28]{Img/T_51c.pdf} \\
		\small $n=5, m=9$ & \small $n=6, m=11$ & \small $n=17, m=29$
	\end{tabular}
	
	\vskip0.2cm
	\fuente{elaboración propia.} 
\end{minipage}
\vspace{3em} % Espacio para separar la minipage del siguiente párrafo

\end{ejemplo}

%Los gráficos anteriores se obtienen con las siguientes instrucciones. 
%\begin{Verbatim}[fontsize=\letraverbatim]
%DiscretePlot3D[8+2*x-x^3-4*y,{x,0,2,0.05},{y,0,1,0.05},
%ExtentSize->Left,BoxRatios->{1,1,1}] \end{Verbatim} 
\vspace{-8mm}
El volumen bajo la figura de $f(x,y)$ sobre el rectángulo $R$ puede verse también como la integral de $A(x),$ 
la cual representa el área de las secciones transversales del sólido;
por tanto el volumen está dado por {\small $\ds \int_a ^b A(x) dx= \int_a ^b \int_c ^d f(x,y) dy dx $.} O de manera similar
{\small $\ds \int_c ^d A(y) dy= \int_c ^d \int_a ^b  f(x,y) dx dy.$}

\vskip0.4cm

Ejemplo de lo anterior es la siguiente integral
que representa el volumen del sólido en mención $$\ds \int_0^2\int_0^1 (8+2x-x^3-4y) dydx = \int_0^1 \int_0^2(8+2x-x^3-4y) dxdy.$$

\begin{figure}[H] 
	\centering
	\caption{Secciones transversales bajo $z = 8+2x - x^3 - 4y$, con $x \in [0,2]$, $y \in [0,1]$.}
	\label{fig:secciones_transversales_bajo}
	\begin{minipage}[b]{0.22\linewidth}
		\centering
		\includegraphics[width=\linewidth]{Img/T_52a.pdf}
	\end{minipage}%
	\begin{minipage}[b]{0.22\linewidth}
		\centering
		\includegraphics[width=\linewidth]{Img/T_52b.pdf}
	\end{minipage}%
	\begin{minipage}[b]{0.22\linewidth}
		\centering
		\includegraphics[width=\linewidth]{Img/T_52c.pdf}
	\end{minipage}%
	\begin{minipage}[b]{0.22\linewidth}
		\centering
		\includegraphics[width=\linewidth]{Img/T_52d.pdf}
	\end{minipage}
	\fuente{elaboración propia.} 
\end{figure}

\subsection*{Integrales triples} \index{Integrales!triples}

Consideraremos una caja rectangular cerrada y acotada $E$ en el espacio con
caras paralelas a los planos coordenados $E=[a,b]\times \lbrack c,d]\times
\lbrack r,s];$ es decir, 
\begin{equation*}
E=\{(x,y,z) \in \mathbb{R}^{3}\,|\,a\leq x\leq b,c\leq y\leq
d,r\leq z\leq s\}.
\end{equation*}
Si $f(x,y,z)$ es una funci\'{o}n continua definida sobre la caja rectangular $E,$
subdividiremos a $E$ en $lnm$ cajas rectangulares m\'{a}s peque\~{n}as a
partir de planos paralelos a los planos coordenados. Esto es, para el
intervalo $[a,b],$ consideremos 
\begin{equation*}
x_{0}=a<x_{1}<x_{2}<\cdots <x_{i}<\cdots <x_{n-1}<x_{n}=b,
\end{equation*}
para el intervalo $[c,d],$ 
\begin{equation*}
y_{0}=c<y_{1}<y_{2}<\cdots <y_{j}<\cdots <y_{m-1}<y_{m}=d,
\end{equation*}
y para el intervalo $[r,s],$ consideremos 
\begin{equation*}
z_{0}=r<z_{1}<z_{2}<\cdots <z_{k}<\cdots <y_{l-1}<y_{l}=s.
\end{equation*}

Si $n,$ $l$ y $m$ crecen entonces el ancho, el largo y el alto de estas
cajas rectangulares decrecen.
Denotaremos $\Delta x_{i}$ al ancho del $i$-\'{e}simo subintervalo sobre el
eje $x,$ $\Delta y_{j}$ es el largo $j$-\'{e}simo subintervalo sobre el eje $y$, y $\Delta z_{k}$ es el alto $k$-\'{e}simo subintervalo sobre el eje $z$, es decir $\Delta x_{i}=x_{i+1}-x_{i},\Delta y_{j}=x_{j+1}-x_{j},$ $\Delta
z_{k}=z_{k+1}-z_{k}.$ 
El volumen de cada una de estas cajas de ancho $\Delta x_{i}$, largo $\Delta y_{j}$ y alto $\Delta z_{k}$ es 
$\Delta V_{ijk}=\Delta x_{i}\Delta y_{j}\Delta z_{k}$. 
La norma de la partici\'{o}n se denota por $\left\Vert P\right\Vert $ y es
el m\'{a}ximo sobre todos los valores $i,$ $j$ y $k$ de todas las
dimensiones de las cajas en que se divide la regi\'{o}n $E.$ 
\vskip0.2cm
Una suma de Riemann de $f(x,y,z) $ sobre la caja rectangular $E$ se consigue al considerar un punto cualquiera 
$\left( x_{i}^{\ast},y_{j}^{\ast },z_{k}^{\ast }\right) ,$ en la $ijk$-\'{e}sima caja
rectangular, esto es $\left( x_{i}^{\ast },y_{j}^{\ast },z_{k}^{\ast}\right) \in \lbrack x_{i-1},x_{i}]\times \lbrack y_{j-1},y_{j}]\times \lbrack z_{k-1},z_{k}],$ lo cual equivale a que $x_{i-1}<x_{i}^{\ast }<x_{i},$ 
$y_{j-1}<y_{j}^{\ast }<y_{j}$ y $z_{k-1}<z_{k}^{\ast }<z_{k}.$ El valor de la
funci\'{o}n es este punto es $f\left( x_{i}^{\ast },y_{j}^{\ast},z_{k}^{\ast }\right) $ 
y este valor se multiplica por el volumen de la $ijk$-\'{e}sima caja rectangular; se suman todos estos productos para obtener
una expresi\'{o}n de la forma 
\begin{equation*}
S_{nml}=\sum_{i=1}^{n}\sum_{j=1}^{m}\sum_{k=1}^{l}f\left( x_{i}^{\ast},y_{j}^{\ast },z_{k}^{\ast }\right) \Delta x_{i}\Delta y_{j}\Delta z_{k}.
\end{equation*}
Dependiendo de la elecci\'{o}n de $\left( x_{i}^{\ast },y_{j}^{\ast },z_{k}^{\ast }\right)$
es posible obtener diversos valores de $S_{nml}.$ 
De manera similar a como se defini\'{o} la integral doble, si el l\'{\i}mite de
las sumas $S_{nml}$ existe para toda elecci\'{o}n que se haga de 
$\left(x_{i}^{\ast },y_{j}^{\ast },z_{k}^{\ast }\right),$ entonces se dice que la
funci\'{o}n $f(x,y,z) $ es integrable sobre la caja rectangular $E,$ y este valor se denota por 
$\ds\iiint\limits_{E}f(x,y,z) dV$ o tambi\'{e}n $\ds \iiint\limits_{E}f(x,y,z) dzdydx.$ Lo anterior nos
motiva a precisar este concepto.

\begin{defi}
La integral triple de una función continua $f(x,y,z)$ 
sobre la caja rectangular $E$ se define como 
\begin{eqnarray}
\lim_{\left\Vert P\right\Vert \rightarrow 0}S_{nml} &=& \lim_{l,m,n \to\infty}\left(\sum_{i=1}^{n}\sum_{j=1}^{m}\sum_{k=1}^{l}f\left( x_{i}^{\ast },y_{j}^{\ast },z_{k}^{\ast}\right) \Delta x_{i}\Delta y_{j}\Delta z_{k} \right)\nonumber \\ &=&\iiint\limits_{E}f\left(x,y,z\right) dV 
\end{eqnarray}
cuando este l\'{\i}mite existe.
\end{defi}
Ilustremos lo anterior con dos integrales triples.
\begin{ejemplo}
Evaluar {\small $\ds \int_{-1}^1\int_0^1\int_0^1  yz^2 e^{x} dzdydx$} mediante sumas de Riemann.
\vskip0.2cm
Al hacer una partición de la región de integración en $nml$ cajas rectangulares se sigue que 
$\Delta x= \frac{2}{n}$ $\Delta y= \frac{1}{m}$ y $\Delta z= \frac{1}{l}$. Además
$x_i= -1+\frac{2i}{n},$ $y_j=\frac{j}{m}$ y $z_k=\frac{k}{l}.$
Por lo tanto 
$$f(x_i,y_j,z_k)\Delta x \Delta y\Delta z=\frac{2 j k^2 e^{\frac{2 i}{n}-1}}{l^3 m^2 n}$$
Con lo cual 
$$\sum_{i=1}^n\sum_{j=1}^m\sum_{k=1}^l f(x_i,y_j,z_k)\Delta x \Delta y\Delta z=\frac{\left(e^2-1\right) (l+1)(2l+1)(m+1)}{6nml^2 \left(1-e^{-2/n}\right)e}$$ entonces 
\begin{multline*} \int_{-1}^1\int_0^1\int_0^1  yz^2 e^{x} dzdydx= \lim_{n,m,l \to \infty}\sum_{i=1}^n\sum_{j=1}^m\sum_{k=1}^l f(x_i,y_j,z_k)\Delta x \Delta y\Delta z \\ =\frac{1}{6} \left(e-\frac{1}{e}\right)\end{multline*}
\end{ejemplo}

\begin{comment}
Con las siguientes instrucciones se evalúan los cálculos anteriores.
 \begin{Verbatim}[fontsize=\letraverbatim]
f[x_,y_,z_]:=z^2*y*Exp[x]
{xi,yj,zk}={-1+2*i/n,j/m,k/l};
{dx,dy,dz}={2/n,1/m,1/l};
suma=Sum[f[xi,yj,zk]*dx*dy*dz,{i,1,n},{j,1,m},{k,1,l}]
Limit[suma,{n,m,l}->{Infinity,Infinity,Infinity}]
Integrate[f[x,y,z],{x,-1,1},{y,0,1},{z,0,1}]
\end{Verbatim}

\end{comment}

\begin{ejemplo}
\begin{sloppypar}
Para $f(x,y,z)= \cos (x) \sin (y)+z^2 \cos^2 (2 y)$ y $E$ el \mbox{paralelepípedo} determinado por 
$[-\frac{\pi}{2},\frac{\pi}{2}]\times [0,\frac{\pi}{4}]\times[0,1]$; evaluar 
{\footnotesize $\ds \iiint_Ef(x,y,z) dzdydx.$}
\vskip0.2cm
Al hacer una partición de $E$ en $nml$ cajas rectangulares es posible establecer que
$\Delta x= \frac{\frac{\pi}{2}-\left(\frac{-\pi}{2}\right)}{n}=\frac{\pi}{2n}$
$\Delta y= \frac{\frac{\pi}{4}-0}{m}=\frac{\pi}{4m}$ y $\Delta z= \frac{1}{l}$. Además
$x_i= -\frac{\pi}{2}+\frac{\pi i}{2n},$ $y_j=\frac{\pi j}{4m}$ y $z_k=\frac{k}{l}.$
Por lo tanto: 
$$f(x_i,y_j,z_k)=\sin\left(\frac{\pi i}{n}\right)\sin\left(\frac{\pi j}{4m}\right)+\frac{k^2}{l^2}\cos^2\left(\frac{\pi j}{2m}\right).$$
Con lo anterior, se establece que $S_{nml}=\sum\limits_{i=1}^n\sum\limits_{j=1}^m\sum\limits_{k=1}^l f(x_i,y_j,z_k)\Delta x \Delta y\Delta z,$ es
{\small 
$$\frac{\pi ^2}{48 m}\left(\frac{(l+1) (2 l+1) (m-1)}{l^2}+\frac{12 \sin \left(\frac{\pi}{8}\right) \sin \left(\frac{\pi  (m+1)}{8 m}\right) \csc \left(\frac{\pi }{8 m}\right)\cot \left(\frac{\pi }{2 n}\right)}{n}\right).$$
}
De modo que $\lim_{n,m,l \to \infty}S_{nml}$ es
\begin{multline*}
\int _{-\frac{\pi }{2}}^{\frac{\pi }{2}}\int _0^{\frac{\pi }{4}}\int _0^1\left(\cos(x)\sin (y)+z^2 \cos ^2(2 y)\right)dzdydx  =2-\sqrt{2}+\frac{\pi ^2}{24}
\end{multline*}
\end{sloppypar}
\end{ejemplo}

\begin{comment}
 \begin{Verbatim}[fontsize=\letraverbatim]
f[x_,y_,z_]:=Cos[x]*Sin[y]+z^2*Cos[2*y]^2
{xi,yj,zk}={-Pi/2+Pi*i/n,Pi/4*j/m,k/l};
{dx,dy,dz}={Pi/n,Pi/(4*m),1/l};
suma=Simplify[
Sum[f[xi,yj,zk]*dx*dy*dz,{i,1,n},{j,1,m},{k,1,l}]];
Limit[suma,{n,m,l}->{Infinity,Infinity,Infinity}]
Integrate[f[x,y,z],{x,-Pi/2,Pi/2},{y,0,Pi/4},{z,0,1}]
\end{Verbatim} 
\end{comment}
Utilizando \verb"Mathematica" es posible simplificar y evaluar las integrales anteriores. 
Su uso se explicará en detalle en los próximos capítulos.
\vskip0.1cm 
%

Para una región acotada sólida $E$ en el espacio, se sigue un proceso análogo al utilizado en las 
integrales dobles. Nótese que $z$ recorre los puntos entre las superficies $z=u_1(x,y)$ y
$z=u_2(x,y)$. La proyección de $E$ en el plano $xy$ muestra que $y$ recorre los puntos que están incluidos entre las curvas $y=h_1(x)$ y $y=h_2(x),$ con $x\in [a,b].$ 


\vskip0.2cm
\noindent
\begin{minipage}{\textwidth}
	\centering
	\captionof{figure}{Integración en el espacio}
	\begin{minipage}{0.39\linewidth}
		\centering
		\includegraphics[width=\linewidth]{Fig/9.pdf}
	\end{minipage}
	
	\vspace{0.3em}
	\parbox{0.95\linewidth}{\centering\fontsize{8pt}{10pt}\selectfont\textbf{Fuente:} elaboración propia.}
\end{minipage}
\vskip0.2cm


Es por esto que la integral respectiva se escribe en la forma
$$\int_a^b\int_{h_1(x)}^{h_2(x)} \int_{u_1(x,y)} ^{u_2(x,y)} f(x,y,z) dz dydx.$$ 

Se debe observar que en algunos casos la proyección del sólido en el plano $xy$ puede corresponder a regiones del tipo I,
tipo II, las cuales fueron descritas anteriormente, o también combinaciones de estas.
\vskip0.2cm
En estos apuntes se consideran otras posibilidades e intercambios en el orden de integración, que dependiendo de la naturaleza de cada problema harán más o menos complicado el tratamiento a estas integrales.

\subsection*{Integrales múltiples sobre regiones parametrizadas}

En \cite{acosta_gempeler_doce_2020}, se considera una función continua $f(x,y,z)$ y una superficie $\sigma$
parametrizada en la forma ${\bf r}(u,v)= \langle x(u,v), y(u,v), z(u,v) \ra,$
con $u \in [a,b]$ y $v \in [c,d];$ cuyas funciones componentes poseen derivadas parciales continuas en 
$[a,b] \times [c,d].$ Al realizar una partición de $\sigma$ en la forma
$$\sigma_{ij}= \{ (x(u_i^\ast,v_j^\ast),y(u_i^\ast,v_j^\ast),z(u_i^\ast,v_j^\ast) )\,|\, u_{i-1} \leq u_i^\ast \leq u_{i}, v_{j-1} \leq v_{j}^\ast \leq v_{j}\} ;$$
considerando el punto $P_{ij}(x(u_i^\ast,v_j^\ast),y(u_i^\ast,v_j^\ast),z(u_i^\ast,v_j^\ast) )$ en cada trozo $\sigma_{ij}$
de la superficie, se obtiene la suma de Riemann de $f$ sobre la superficie $\sigma$ parametrizada por ${\bf r}(u,v)$ y con la partición de $\sigma$,
$$ S_{nm}=\sum_{i=1}^n \sum_{j=1}^m f(x(u_i^\ast,v_j^\ast),y(u_i^\ast,v_j^\ast),z(u_i^\ast,v_j^\ast) ) \Vert {\bf r}_u \times {\bf r}_v \Vert \Delta u_i \Delta v_j .$$
El término $\Vert {\bf r}_u \times {\bf r}_v \Vert $ es el área del paralelogramo determinado por
los vectores ${\bf r}_u$ y ${\bf r}_v.$
Con un procedimiento como el utilizado cuando presentamos integrales dobles, se sigue que en este caso la integral de $f$ 
sobre la superficie $\sigma$ con la parametrización ${\bf r}(u,v)$ está dada por
\begin{multline}
\iint_\sigma f dA = \lim_{n,m \to \infty} S_{nm}  = \int_a^b\int_c^d f({\bf r}(u,v)) \Vert {\bf r}_u \times {\bf r}_v \Vert dudv  \label{integralesdobles2}
\end{multline}
Si la parametrización ${\bf r}(u,v,w)= \langle x(u,v,w), y(u,v,w), z(u,v,w) \ra,$
de un sólido $E$, con $u \in [a,b],$ $v \in [c,d]$ y $w \in [p,q]$;
con funciones componentes que poseen derivadas parciales continuas en 
$[a,b] \times [c,d] \times [p,q].$

\begin{figure}[H]
	\centering
	\caption{Transformación de una región en el espacio}
	\label{fig:transform-region}
	\begin{minipage}[b]{0.30\linewidth}
		\centering
		\includegraphics[width=\linewidth]{Fig/10.pdf}
	\end{minipage}
	\begin{minipage}[b]{0.30\linewidth}
		\centering
		\includegraphics[width=\linewidth]{Fig/11.pdf}
	\end{minipage}
	\begin{minipage}[b]{0.30\linewidth}
		\centering
		\includegraphics[width=\linewidth]{Fig/12.pdf}
	\end{minipage}
	\parbox{0.95\linewidth}{\centering\fontsize{9pt}{10pt}\selectfont\textbf{Fuente:} elaboración propia.}
\end{figure}
\vskip0.2cm

Es posible realizar una partición de $E$ en la forma
\[E_{ijk}= \{ P_{ijk}^\ast\,|\, u_{i-1} \leq u_i^\ast \leq u_{i}, v_{j-1} \leq v_{j}^\ast \leq v_{j}, w_{k-1} \leq w_{k}^\ast \leq w_{k}\};\]
En donde $P_{ijk}^\ast=(x(u_i^\ast,v_j^\ast,w_k^\ast),y(u_i^\ast,v_j^\ast,w_k^\ast),z(u_i^\ast,v_j^\ast,w_k^\ast) );$
al considerar el punto $P_{ijk}(x(u_i^\ast,v_j^\ast,w_k^\ast),y(u_i^\ast,v_j^\ast,w_k^\ast),z(u_i^\ast,v_j^\ast,w_k^\ast) )$
en cada trozo $E_{ijk}$ del sólido, se obtiene la suma de Riemann de $f$ sobre el sólido $E$ parametrizado por 
${\bf r}(u,v,w)$ y con la partición de $E$,
$$ S_{nml}=\sum_{i=1}^l \sum_{i=1}^n \sum_{j=1}^m f(P_{ijk}^\ast) | {\bf r}_u \cdot ({\bf r}_v \times {\bf r}_w) | \Delta u_i \Delta v_j \Delta w_k.$$
El término $| {\bf r}_u \cdot ({\bf r}_v \times {\bf r}_w) |$ es el volumen del paralelepípedo determinado por
los vectores ${\bf r}_u,$ ${\bf r}_v$ y ${\bf r}_w$. 
De manera análoga al procedimiento utilizado cuando presentamos integrales triples, se sigue que la integral de $f$ 
sobre $E$ con la parametrización ${\bf r}(u,v,w)$ está dada por
\begin{multline} \label{math1}
\ds \iiint_E f dV = \lim_{l,m,n \to \infty} S_{nml}  \\ 
=\int_a^b\int_c^d \int_p^q f({\bf r}(u,v,w)) | {\bf r}_u \cdot ({\bf r}_v \times {\bf r}_w) | dudvdw 
\end{multline}
Al término $|{\bf r}_u \cdot ({\bf r}_v \times {\bf r}_w) |dudvdw$ lo denominaremos elemento diferencial de volumen.
\index{Elemento diferencial!de volumen}
Esta última expresión puede verse como la integral de la función $f$ si el sólido $E^*$ es
parametrizado por la función ${\bf r}(u,v,w).$ Si en la expresión (\ref{math1}) $f(x,y,z)=1$ esta integral determina 
el volumen $V$ del sólido $E$, es decir
\begin{eqnarray} \label{math2}
V=\iiint_{E}|{\bf r}_u \cdot ({\bf r}_v \times {\bf r}_w) | dudvdw
\end{eqnarray}

\subsection*{Aplicaciones de las integrales triples}

De manera similar al caso bidimensional la masa $M,$ de un sólido $E$ con densidad $\rho (x,y,z)$ se determina por 
$$M=\ds \iiint_E \rho (x,y,z) dV$$ y los momentos con respecto a los planos coordenados se determinan por 
$$\ds \mu_{xy}=\iiint_E z\rho(x,y,z) dV,\,\,\, \mu_{xz}=\iiint_E y\rho(x,y,z) dV,$$
$$\ds  \mu_{yz}=\iiint_E x\rho(x,y,z)dV.$$
De tal forma que el centro de masa de $E$ es
$(\overline{x},\overline{y},\overline{z})= (\frac{\mu_{yz}}{M}, \frac{\mu_{xz}}{M},\frac{\mu_{xy}}{M}) .$
Una exposición más detalladas de estos elementos y otros usos de las integrales múltiples puede consultarse en \cite{thomas_calculo_2015}.
\index{Centro de masa}
\subsection*{Integrales de superficie}
Una parametrizaci\'{o}n de una superficie es una funci\'{o}n 
$\phi:D\subseteq \mathbb{R}^{2}\rightarrow \mathbb{R}^{3}.$
La superficie $S$ correspondiente a la
funci\'{o}n $\phi $ es su imagen $\phi \left( D\right) =S$ se escribe 
$\phi(u,v)=\left\langle x\left( u,v\right) ,y\left( u,v\right),z\left( u,v\right) \right\rangle .$
La integral de un campo escalar $f$ sobre una superficie $S,$  está dada por

\begin{equation} \label{areasup1}
\iint_{S}f(x,y,z) dS=\iint_{D}f\left( \phi \left( u,v\right)
\right) \left\| {\bf T}_{u}\times {\bf T}_{v}\right\| dudv
\end{equation}

siendo
$${\bf T}_{u}=\frac{\partial \phi }{\partial u}=\frac{\partial x}{\partial u}\vi+\frac{\partial x}{\partial u}\vj+\frac{\partial x}{\partial u}\vk$$
y $${\bf T}_{v}=\frac{\partial \phi }{\partial v}=\frac{\partial x}{\partial v}\vi+\frac{\partial x}{\partial v}\vj+\frac{\partial x}{\partial v}\vk.$$

La integraci\'{o}n de una campo vectorial ${\bf f}$ sobre $S,$ está dada por
\begin{equation}
\iint_{S} {\bf f}(x,y,z) \cdot dS=\iint_{D} {\bf f} \left( \phi \left(
u,v\right) \right) \cdot \left( {\bf T}_{u}\times {\bf T}_{v}\right) dudv
\end{equation}
Una exposición más detallada de esta formulación puede encontrarse en \cite{marsden_calculo_2013}.
\vskip0.2cm
Si en la expresión (\ref{areasup1}) $f(x,y,z)=1,$ esta integral se reduce al área $A$ de la superficie $D,$ es decir
\begin{equation} \label{areasup}
A=\iint_{D} \left\| {\bf T}_{u}\times {\bf T}_{v}\right\| dudv
\end{equation}
Al término $\left\| {\bf T}_{u}\times {\bf T}_{v}\right\|dudv$ lo denominaremos elemento diferencial de área.
\index{Elemento diferencial!de área}

\subsection*{Cambio de variables en integrales dobles} \label{cambioidobles}
\vskip0.2cm
Con el propósito de evaluar la integral $\ds \iint\limits_{R}f(x,y) dA$
se hace el cambio $(x,y) ={\bf r}(u,v)$ para transformar la región $R$ en el sistema 
$(x,y) $ en una región $R^{\ast }$ en el sistema coordenado  $(u,v).$
Un rectángulo en el sistema $uv$ con lados de longitud $\Delta u$ de largo y $\Delta v$ de ancho se transforma mediante la función ${\bf r}.$
\vskip0.2cm
Asumiendo que ${\bf r}$ es diferenciable, si $(u,v) $ son las coordenadas del rectángulo, los v\'{e}rtices
adyacentes son $( u+\Delta u,v) $ y $(u,v+\Delta v)$; mediante la aplicación ${\bf r},$ el rectángulo
se transforma en una región cuya área se puede aproximar mediante el área de un paralelogramo. Además, 
aplicando el teorema del valor medio, se sigue que

\begin{eqnarray*}
{\bf r}(u+\Delta u,v) -{\bf r}(u,v) &=& \frac{{\bf r}( u+\Delta u,v)-{\bf r}(u,v)}{\Delta u}\Delta u 
\approx \frac{\partial {\bf r}}{\partial u}\left( u^{\ast},v\right)\Delta u \\
{\bf r}(u,v+\Delta v) -{\bf r}(u,v) &=& \frac{{\bf r}(u,v+\Delta v)-{\bf r}(u,v)}{\Delta v}\Delta v 
\approx \frac{\partial{\bf r}}{\partial u}( u,v^{\ast})\Delta v.
\end{eqnarray*}

El área de esta pequeña región est\'{a} dada por 
\begin{equation}
\Delta A=\left\Vert \frac{\partial {\bf r}}{\partial u}\Delta u\times 
\frac{\partial {\bf r}}{\partial v}\Delta v\right\Vert =\left\Vert \frac{\partial {\bf r}}{\partial u}\times \frac{\partial {\bf r}}{\partial v}\right\Vert \Delta u\Delta v
\end{equation}
que se puede abreviar en la forma
$\Delta A=\left\vert \frac{\partial (x,y)}{\partial (u,v)}\right \vert \Delta u\Delta v.$
Por simplicidad $x=x(u,v)$ y $y=y(u,v)$, entonces
\begin{equation}
\left\vert \frac{\partial (x,y) }{\partial(u,v)}\right\vert =\left\Vert\frac{\partial{\bf r}}{\partial u}\times \frac{\partial{\bf r}}{\partial v}\right\Vert =
\left\vert \begin{array}{cc}\frac{\partial x}{\partial u}&\frac{\partial x}{\partial v}\\[.1cm]  
\frac{\partial y}{\partial u} & \frac{\partial y}{\partial v}
\end{array} \right\vert
\end{equation}

En algunos libros como \cite{stewart_calculo_2018}, el t\'{e}rmino $\left\vert \frac{\partial\left(x,y\right)}{\partial(u,v)}\right\vert$ 
se denomina el jacobiano de la transformación ${\bf r}$. \index{Jacobiano}
Con los elementos anteriores se establece la siguiente igualdad.
\begin{equation}
\iint\limits_{R}f(x,y) dA=\iint\limits_{R^{\ast }}f\left({\bf r}(u,v) \right) \left\vert \frac{\partial (x,y) }{\partial (u,v)}\right\vert dudv \end{equation}

\subsection*{Cambio de variables en integrales triples}
Con el propósito de evaluar la integral 
$$\iiint\limits_{R}f(x,y,z) dV$$
se puede hacer el cambio $(x,y,z) ={\bf r}(u,v,w)$ para transformar la región $R$ en el sistema 
de coordenadas $(x,y,z) $ en una region $R^{\ast }$ en el sistema coordenado  $(u,v,w) .$
Consideremos un cubo en el sistema $uvw,$ con aristas de longitud $\Delta u$ de largo, $\Delta v$ de ancho y $\Delta w$ de altura, que se transforma mediante la función ${\bf r}$. \vskip0.2cm 
Asumiendo que ${\bf r}$ es diferenciable,
si $(u,v,w) $ son las coordenadas de un vértice del cubo, entonces los v\'{e}rtices
adyacentes son $\left( u+\Delta u,v,w\right) ,$ $\left( u,v+\Delta v,w\right)$ y $\left( u,v,w+\Delta w\right) ;$ mediante la aplicación ${\bf r},$ el cubo se transforma en una región cuyo volumen se puede aproximar mediante un \mbox{paralelep\'{\i}pedo}.
Aplicando el teorema del valor medio, se sigue que

\begin{eqnarray*}
{\bf r}\left( u+\Delta u,v,w\right) -{\bf r}(u,v,w) &=& \frac{{\bf r}\left( u+\Delta u,v,w\right) -{\bf r}(u,v,w) }{\Delta u}\Delta u \\[0.1cm] 
&\approx &   \frac{\partial {\bf r}}{\partial u}\left( u^{\ast},v,w\right) \Delta u \\
{\bf r}\left( u,v+\Delta v,w\right) -{\bf r}(u,v,w) &=& \frac{{\bf r}(u,v+\Delta v,w)-{\bf r}(u,v,w)}{\Delta v}\Delta v \\[0.1cm]
&\approx& \frac{\partial{\bf r}}{\partial v}\left( u,v^{\ast},w\right)\Delta v\\
{\bf r}\left( u,v,w+\Delta w\right) -{\bf r}(u,v,w) &=& \frac{{\bf r}\left(
u,v,w+\Delta w\right) -{\bf r}(u,v,w)}{\Delta w}\Delta w \\[0.1cm]
&\approx & \frac{\partial {\bf r}}{\partial w}\left( u,v,w^{\ast}\right) \Delta w,
\end{eqnarray*}

El volumen de esta región, $\Delta V$, est\'{a} dada por 
\begin{equation}
\left\vert \left( \frac{\partial {\bf r}}{\partial u}\Delta u\times 
\frac{\partial {\bf r}}{\partial v}\Delta v\right) \cdot \frac{\partial {\bf r}}{\partial w}\Delta w\right\vert =\left\vert \left( \frac{\partial {\bf r}}{\partial u}\times \frac{\partial {\bf r}}{\partial v}\right) \cdot \frac{\partial {\bf r}}{\partial w}\right\vert \Delta u\Delta v\Delta w
\end{equation}
que se puede abreviar en la forma
$\Delta V=\left\vert \frac{\partial (x,y,z) }{\partial \left(
u,v,w\right) }\right\vert \Delta u\Delta v\Delta w.$
Por simplicidad, si $x=x(u,v,w) ,$ $y=y(u,v,w) $ y $z=z(u,v,w) $ entonces
\begin{equation}
\left\vert \frac{\partial (x,y,z) }{\partial(u,v,w)}\right\vert =\left\vert\left(\frac{\partial{\bf r}}{\partial u}\times \frac{\partial{\bf r}}{\partial v}\right) \cdot \frac{\partial{\bf r}}{\partial v}\right\vert =
\left\vert 
\begin{array}{ccc}\frac{\partial x}{\partial u}&\frac{\partial x}{\partial v}&\frac{\partial x}{\partial w}\\[.1cm]  
\frac{\partial y}{\partial u}&\frac{\partial y}{\partial v}&\frac{\partial y}{\partial w} \\[0.1cm] 
\frac{\partial z}{\partial u}&\frac{\partial z}{\partial v}&\frac{\partial z}{\partial w}\end{array}\right\vert 
\end{equation}

En algunos libros el t\'{e}rmino $\left\vert \frac{\partial\left(x,y,z\right)}{\partial(u,v,w)}\right\vert$ 
se denomina el jacobiano de la transformación ${\bf r}$. \index{Jacobiano}
Esto muestra la siguiente igualdad.
\begin{equation}
\iiint\limits_{R}f(x,y,z) dV=\iiint\limits_{R^{\ast}}f\left({\bf r}(u,v,w) \right)\left\vert\frac{\partial (x,y,z)}{\partial (u,v,w) }\right\vert dudvdw
\end{equation}

Los siguientes casos particulares de cambio de variables en integrales dobles y triples son utilizados 
de manera frecuente.
\subsection*{Coordenadas polares, cilíndricas y esféricas} \label{esfericas}
\index{Coordenadas!polares} \index{Coordenadas!cilíndricas} \index{Coordenadas!esféricas}
Para una circunferencia de radio $r,$ la longitud $s$, de un arco de $\theta$ radianes es $r\theta$.
El área de un sector circular de arco $\Delta \theta$ y radio $r$ está dado por 
$\frac{1}{2}r^2 \Delta \theta$ y el área $\Delta A,$ de la región acotada por dos sectores 
circulares de radio $r_i$ y $r_{i+1}$ (con $r_i<r_{i+1}$) está dada por
\begin{eqnarray*}
\Delta A &=& \frac{1}{2}r_{i+1}^2 \Delta \theta-\frac{1}{2}r_{i}^2 \Delta \theta 
= \frac{r_{i+1}+r_i}{2}(r_{i+1}-r_i)\Delta \theta  = r_i^*\Delta r \Delta \theta 
\end{eqnarray*}

%\begin{figure}
    %\centering
    %\includegraphics[width=0.5\linewidth]{Fig/12.pdf}
    %\caption{Transformación a coordenadas polares}
    %\label{fig:polar-transform}
%\end{figure}
\noindent
\begin{minipage}{\textwidth}
	\centering
	\captionof{figure}{Transformación a coordenadas polares}
	\label{fig:falta_título}
	\begin{minipage}{0.28\linewidth}
		\centering
		\includegraphics[width=\linewidth]{Fig/13.pdf}
	\end{minipage}
	\hspace{1em}
	\begin{minipage}{0.28\linewidth}
		\centering
		\includegraphics[width=\linewidth]{Fig/14.pdf}
	\end{minipage}
	
	\vspace{0.3em}
	\parbox{0.95\linewidth}{\centering\fontsize{9pt}{10pt}\selectfont\textbf{Fuente:} elaboración propia.}
\end{minipage}




Teniendo en cuenta las relaciones anteriores y el teorema de cambio de variables se puede establecer que en el sistema de coordenadas polares, $x=r \cos \theta$, $y=r \sin \theta$,  la transformación de una integral doble es
{\small
\begin{equation*}\label{integraldoblepolares}
\iint_R f(x,y) dydx = \iint_S f(r \cos \theta,r \sin \theta) r drd\theta.
\end{equation*}
}
\vspace{-7.9mm}


Utilizando el sistema de coordenadas cilíndricas $x=r\cos \theta ,$ 
$y=r\sin \theta ,$ $z=z;$ o el sistema de coordenadas esf\'{e}ricas,
$x=\rho \cos \theta \sin \phi ,$ $y=\rho \sin \theta \sin \phi ,$ $z=\rho\cos \phi ,$ 
con los siguientes gráficos se ilustra la obtención del elemento diferencial de volumen en estos casos.


\vskip0.2cm
\noindent
\begin{minipage}{\textwidth}
	\centering
	\captionof{figure}{Vectores ${\bf F}(x,y,z),$ ${\bf T}$ y ${\bf n}$}
	\begin{minipage}{0.31\linewidth}
		\centering
		\includegraphics[width=\linewidth]{Fig/15.pdf}
	\end{minipage}
	
	\vspace{0.3em}
	\parbox{0.95\linewidth}{\centering\fontsize{8pt}{10pt}\selectfont\textbf{Fuente:} elaboración propia.}
\end{minipage}
\vskip0.2cm


El jacobiano de la transformación a coordenadas cilíndricas se obtiene a partir de la relación
$$\left\vert \begin{array}{ccc}
\cos \theta  & \sin \theta  & 0 \\ 
-r\sin \theta  & r\cos \theta  & 0 \\ 
0 & 0 & 1\end{array}\right\vert = r\cos ^{2}\theta +r\sin ^{2}\theta = r,$$
y el elemento diferencial de Volumen es $dV= r \, dr d\theta dz,$ 
por lo tanto, la fórmula \ref{math1} se escribe como
\index{Integrales triples!en coordenadas cilíndricas} 
$$\iiint_{E}f(x,y,z)dV=\iiint_{E^{\ast }}f(r\cos \theta ,r\sin \theta,z)r\, dzdrd\theta .$$


\vspace{-6mm}
\vskip0.2cm
\noindent
\begin{minipage}{\textwidth}
	\centering
	\vskip0.2cm
	\captionof{figure}{Elemento diferencial de volumen en coordenadas esféricas}
	\label{fig:volumen-esfericas}
	\begin{minipage}{0.31\linewidth}
		\centering
		\includegraphics[width=\linewidth]{Fig/16.pdf}
	\end{minipage}%
	\begin{minipage}{0.31\linewidth}
		\centering
		\includegraphics[width=\linewidth]{Fig/17.pdf}
		\label{fig:esfericas}
	\end{minipage}
	
	\vspace{0.3em}
	\parbox{0.95\linewidth}{\centering\fontsize{8pt}{10pt}\selectfont\textbf{Fuente:} elaboración propia.}
\end{minipage}
\vskip0.2cm


El jacobiano en el caso de las coordenadas esf\'{e}ricas corresponde a 
$$\left\vert \begin{array}{ccc}
\cos \theta \sin \phi  & \sin \theta \sin \phi  & \cos \phi  \\ 
-\rho \sin \theta \sin \phi  & \rho \cos \theta \sin \phi  & 0 \\ 
\rho \cos \theta \cos \phi  & \rho \sin \theta \cos \phi  & -\rho \sin \phi 
\end{array}\right\vert = -\rho ^{2}\sin \phi,$$
que determina el elemento diferencial de volumen $dV= \rho ^2 \sin \phi d\rho d\theta d\phi$
y permite obtener la siguiente igualdad:
\begin{multline}
\iiint_{E}f(x,y,z)dV \\ =\iiint_{E^{\ast }}f(\rho \cos \theta \sin \phi ,\rho
\sin \theta \sin \phi ,\rho \cos \phi )\rho ^{2}\sin \phi\, d\rho d\theta d\phi 
\end{multline}
En el apéndice {B} el lector puede consultar las generalidades de otros sistemas de coordenadas.
\index{Integrales triples!en coordenadas Esféricas}

\begin{ejemplo}
Si $E$ es el sólido acotado por los paraboloides $z=2-x^2-y^2$ y $z=x^2+y^2.$
Evaluar la integral $\displaystyle \iiint_E y^2 dV.$ La intersección de los paraboloides está dada por $x^2+y^2=1,$ que corresponde a un círculo unitario 
en el plano $xy;$ con estos elementos se plantea y evalúa esta integral:

\vspace{-4mm}

\begin{multline*}
\int_{-1}^{1}\int_{-\sqrt{1-x^{2}}}^{\sqrt{1-x^{2}}}\int_{x^{2}+y^{2}}^{2-x^{2}-y^{2}}y^{2}dzdydx = \int_{-1}^{1}\int_{-\sqrt{1-x^{2}}}^{\sqrt{1-x^{2}}}\bigg[ y^{3}\bigg]_{z=x^{2}+y^{2}}^{z=2-x^{2}-y^{2}} dydx \\
=\int_{-1}^{1}\int_{-\sqrt{1-x^{2}}}^{\sqrt{1-x^{2}}}\left( 2y^{2}\left(1-x^{2}-y^{2}\right) \right) dydx
= \int_{-1}^{1}\frac{8}{15}\left( 1-x^{2}\right) ^{\frac{5}{2}}dx= \frac{\pi}{6}.
\end{multline*}

Aplicando un cambio a coordenadas cilíndricas se obtiene
\begin{multline*}
\int_{0}^{2\pi }\int_{0}^{1}\int_{r^{2}}^{2-r^{2}}r\left( r\sin \theta\right) ^{2}dzdrd\theta \\
= \int_{0}^{2\pi}\int_{0}^{1}2r^{3}\left( 1-r^{2}\right) \sin ^{2}\theta drd\theta 
= \int_{0}^{2\pi }\frac{1}{6}\sin ^{2}\theta d\theta = \frac{\pi}{6} 
\end{multline*}
cuyos valores son iguales, como se había establecido teóricamente.
\end{ejemplo}
\section{Divergencia y rotacional} \label{divergenciayrotacional}
Para un campo vectorial 
${\bf F}(x,y,z)= \langle P(x,y,z),Q(x,y,z),R(x,y,z) \ra,$ definido en $\mathbb{R}^3,$ cuyas funciones componentes poseen derivadas parciales continuas,
se define el rotacional y la divergencia de ${\bf F}$ que se denotan por $\rot{F}$ y ${\rm div\,}{\bf F},$ y 
están dados, respectivamente, por 
\index{Divergencia} \index{Rotacional} 
$$\rot{F}(x,y,z)= \langle  \frac{\partial R}{\partial y}-\frac{\partial Q}{\partial z}, \frac{\partial P}{\partial z}-\frac{\partial R}{\partial x}, \frac{\partial Q}{\partial x}-\frac{\partial P}{\partial y} \rangle $$
y $${\rm div\,}{\bf F}(x,y,z)= \frac{\partial P}{\partial x}+\frac{\partial Q}{\partial y}+\frac{\partial R}{\partial z}.$$
Nótese que $\rot{F}(x,y,z)$ es un campo vectorial y ${\rm div\,}{\bf F}(x,y,z)$ es un campo escalar. 
Existe una manera de recordar brevemente dichos cálculos, en términos del operador $\nabla$, que se presentó en la sección 3.3.
\begin{eqnarray*}
{\rm rot}{\bf F} &=& \nabla \times{\bf F}  = \left\vert \begin{array}{ccc}
\vi & \vj &  \vk \\  \frac{\partial }{\partial x}  & \frac{\partial }{\partial y}
  & \frac{\partial }{\partial y}  \\ P & Q & R \end{array}\right\vert\\
{\rm div}{\bf F} &=& \nabla \cdot{\bf F}.
\end{eqnarray*}

\begin{ejemplo}
Para el campo vectorial ${\bf F}(x,y,z) =\left\langle xye^{z},z\cos (xy) ,\sin (xy) \right\rangle,$
es claro que 
$${\rm div}{\bf F}=\frac{\partial }{\partial x} \left( xye^{z}\right) +\frac{\partial }{\partial y} \left( z\cos (xy) \right) +\frac{\partial }{\partial z}(\sin(xy))= ye^{z}-xz\sin xy,$$ y también
$${\rm rot}{\bf F} = \left\langle (x-1) \cos(xy),y(xe^{z}-\cos(xy)) ,-xe^{z}-yz\sin xy \right\rangle.$$
\end{ejemplo}
La integración de estos tipos especiales de funciones se considerarán en el capítulo \ref{chap:stokes},
así como algunas de sus aplicaciones.
\newpage